\documentclass[11pt]{article}
\usepackage[margin=1.1in]{geometry}
\usepackage[T1]{fontenc}
\usepackage{lmodern}
\usepackage{microtype}
\usepackage{natbib}
\usepackage{../shared/smcbvs}
\bibliographystyle{abbrvnat}
\setcitestyle{authoryear,round}
\linespread{1.08}
\newcommand{\bridgepending}{The short bridging argument that puts the statement in this reader form (set-size
arithmetic or assembly from the checked results) awaits its final independent check.}

\title{What we prove: consistency of the SMC-BVS method}
\author{SMC-BVS project}
\date{Draft, \today}

\begin{document}
\maketitle

\begin{abstract}
\noindent
This section collects every consistency result we have for the SMC-BVS method, a sequential Monte Carlo method for
Bayesian variable selection in the linear model. The results answer four questions, in order. Does the posterior pick
the right model? Is the posterior right about the distribution of the response? Does the method, as coded, deliver that
posterior? And what carries over to the parallel version of the method? Every result is asymptotic in the sample size
$n$, allows many more predictors than observations, and concerns the prior exactly as the software sets it, with no
tuning by the user.
\end{abstract}

% =====================================================================================================
\section{Overview}\label{sec:overview}

\emph{This section lists the results in words and says where each one is stated.}

\paragraph{The posterior picks the right model (Section~\ref{sec:selection}).} For fixed designs with a sparse
eigenvalue floor and signals above $\rn$, the posterior puts mass at least $1-C\rn$ on the true model, on average and
uniformly over such designs; the highest posterior model and the median probability model equal the true model with
probability tending to one, and every inclusion probability is within $\OP(\rn)$ of the truth. The same holds for
random Gaussian designs under a population condition, and, with an omitted-variables twist, when some true predictors
are missing from the data. The rate $\rn$ cannot be improved in general.

\paragraph{The posterior is right about the distribution (Section~\ref{sec:distribution}).} The posterior contracts
around the true conditional law of the response given the predictors, at the rate $\rn$ in the Hellinger distance of
\citet{SongLiang2015}, with no beta-min condition; at the parametric rate when the true predictors are detectable; and
the same holds for the part of the posterior that a completed run returns.

\paragraph{The method, as coded, delivers it (Section~\ref{sec:delivery}).} The method's start law already concentrates
on the true model. On one data event of probability tending to one, every run that produces output reports exactly the
true model, whatever its seed. A certified small model rarely misses a true predictor, under conditions much weaker than
those of the main theorem. Every printed estimate is consistent, and the evidence estimate is unbiased.

\paragraph{The parallel method (Section~\ref{sec:parallel}).} The pooling of the four islands is covered; the rest is
empirical or open, and we say so.

\paragraph{What ``asymptotic'' means here.} Every result is a statement as $n\to\infty$, with $p$ allowed to grow
polynomially in $n$ and to exceed $n$. At moderate $n$ and large $p$ the posterior is not close to a point mass: it is a
mixture over the true model and its many one-predictor supersets, each individually light (Remark~\ref{rem:moderate}).
The theorems describe the limit of this picture.

\paragraph{Status notes.} Each result is followed by a shaded note giving its verification status: ``proved and
independently checked'', ``proved and checked, conditional on'' named inputs, or ``in progress''. Nothing in progress is
called a theorem. The notes are for internal use and are removed for submission; Appendix~\ref{app:status} gives the
underlying records.

\begin{table}[ht]
\centering\small
\caption{The results of this paper.}\label{tab:overview}
\begin{tabularx}{\linewidth}{@{}l>{\raggedright\arraybackslash}X l l@{}}
\toprule
Result & In words & Design & Status \\
\midrule
Thm~\ref{thm:post-fixed} & posterior concentrates on $\Sz$ at rate $\rn$ & fixed & checked$^{a,e}$ \\
Thm~\ref{thm:post-margin} & the same under an identifiability margin & fixed & conditional$^{b,e}$ \\
Thm~\ref{thm:post-random} & the same for Gaussian rows & random & checked$^{e}$ \\
Thm~\ref{thm:sharp} & the rate cannot be improved & both & checked$^{e}$ \\
Thm~\ref{thm:omitted} & omitted variables: support of the best candidate predictor & random & conditional$^{b,e}$ \\
Thm~\ref{thm:dist} & posterior recovers the conditional law at rate $\rn$ & random & checked \\
Thm~\ref{thm:dist-param} & parametric rate for detectable signals & random & checked$^{c}$ \\
Cor~\ref{cor:dist-output} & the same for the list a completed run returns & random & checked$^{d}$ \\
Thm~\ref{thm:start} & the start law concentrates on $\Sz$ & random & checked \\
Thm~\ref{thm:reported-random} & every run that produces output reports $\Sz$ (main theorem) & random & checked$^{e}$ \\
Thm~\ref{thm:reported-fixed} & the same, fixed design & fixed & checked$^{e}$ \\
Thm~\ref{thm:screen}, \ref{thm:screen-margin} & a certified small model rarely misses a truth & fixed & checked \\
Thm~\ref{thm:outputs} & every printed estimate is consistent & random & checked$^{e}$ \\
Rem~\ref{rem:unbiased} & the pooled evidence estimate is unbiased & any & per island checked; pooled not yet \\
Sec~\ref{sec:parallel} & parallel engine; many islands & --- & in progress / open \\
\bottomrule
\end{tabularx}
\par\smallskip\raggedright\footnotesize
$^{a}$ In-probability form; the expectation bound is conditional on two supporting lemmas not yet independently checked.
$^{b}$ Conditional on those two lemmas.
$^{c}$ Its one input, Theorem~\ref{thm:post-fixed}, is checked.
$^{d}$ Under exact arithmetic.
$^{e}$ The short bridging arguments that put the statement in reader form await their final independent check.
\end{table}

% =====================================================================================================
\section{Setting and notation}\label{sec:setting}

\emph{This section fixes the data, the model, the prior and the posterior, and names the few objects of the method that
the theorems talk about. The method itself is described in full in the methodology section; here we recall only what
the statements need. Every symbol is listed in the notation table, Appendix~\ref{app:notation}.}

\paragraph{Notation.}
For an integer $m\ge1$, $\set{m}:=\{1,\dots,m\}$. Vectors are columns; $\norm{v}$ is the Euclidean norm and
$\norm{v}_1$ the $\ell_1$ norm. $I_n$ is the $n\times n$ identity and $\one_n$ (or $\one$) the vector of ones. For a
symmetric matrix $B$, $\lmin(B)$ is its smallest eigenvalue. For a set $S$, $\abs{S}$ is its number of elements.
$\log$ is the natural logarithm and $\log_2$ the binary one; $\lceil\cdot\rceil$ and $\lfloor\cdot\rfloor$ round up and
down; $a\vee b:=\max(a,b)$. $\ind{E}$ is the indicator of an event or condition $E$.

\subsection{Data and model}\label{sec:model}

We observe a response vector $\yr\in\R^n$ and an $n\times\praw$ matrix of predictors $\Xr$. The method first removes
every predictor that is constant in the sample. We write $p$ for the number of \emph{retained} predictors, indexed by
$\set{p}$. Everything the method computes, including its prior, uses these $p$ predictors only.

The data follow a Gaussian linear model with an intercept:
\begin{equation}\label{eq:model}
  \yr=\alpha^*\one_n+\Xr\bstar+\varepsilon,\qquad \varepsilon\sim\Normal_n(0,\sigs^2I_n),
\end{equation}
with $\sigs^2>0$ and $\varepsilon$ independent of $\Xr$. In words: each response is a linear function of the predictors
plus independent Gaussian noise of variance $\sigs^2$. The \emph{true model} is the set of predictors with a nonzero
coefficient, $\Sz:=\supp\bstar=\{j:\beta^*_j\neq0\}$, and $\sz:=\abs{\Sz}$ is its size.

We consider two kinds of design.
\begin{itemize}[leftmargin=1.5em]
  \item \emph{Fixed design}: $\Xr$ is deterministic, and probability refers to the noise $\varepsilon$ only.
  \item \emph{Random Gaussian design}: the rows of $\Xr$ are independent and identically distributed (i.i.d.)
    $\Normal_p(m_n,\Sig)$. Because the method standardises every column (Section~\ref{sec:standardise}), we may and do take
    $\Sig$ to be a correlation matrix (unit diagonal). Probability then refers to the joint law of $(\Xr,\varepsilon)$.
    We write $W=(W_1,\dots,W_p)\sim\Normal_p(0,\Sig)$ for a generic population row; for $S\subseteq\set{p}$, $W_S$ is the
    sub-vector indexed by $S$ and $\Sigma_{n,S}$ the principal sub-matrix of $\Sig$ indexed by $S$.
\end{itemize}
A law for the rows is an assumption about how the data are generated, not a choice of a particular family of designs.

\subsection{What the method sees: standardised data}\label{sec:standardise}

Let $P^\perp_{\one}:=I_n-n^{-1}\one\one^\top$ be the centring matrix and $D$ the diagonal matrix of the sample standard
deviations of the retained columns of $\Xr$. The method works with the centred and scaled data
\begin{equation}\label{eq:standardise}
  X:=P^\perp_{\one}\Xr D^{-1},\qquad
  y:=P^\perp_{\one}\yr/s_y,\qquad s_y^2:=\norm{P^\perp_{\one}\yr}^2/(n-1).
\end{equation}
Each column $x_j$ of $X$ has mean zero and $\norm{x_j}^2=n-1$, and $y$ has mean zero and sample variance one. On this
scale the true coefficients and the mean signal are
\[
  \tstar:=D\bstar\quad\text{(standardised coefficients)},\qquad \mu:=P^\perp_{\one}\Xr\bstar=X\tstar\in\R^n,
\]
and $\tmin:=\min_{j\in\Sz}\abs{\theta^*_j}$ is the smallest true standardised coefficient. Standardising does not move
the support: $\supp\tstar=\Sz$. For a fixed design $\tstar$ is deterministic; for a random design it is random through
$D$, and a theorem that places its conditions on $\tstar$ says so.

For $S\subseteq\set{p}$: $X_S$ is the $n\times\abs{S}$ matrix of the columns in $S$; $G_S:=X_S^\top X_S$ is their Gram
matrix; $\Rhat_S:=G_S/(n-1)$ is their sample correlation matrix; and $P_S$ is the orthogonal projection onto the span of
the columns of $X_S$ ($P_\emptyset:=0$).

\subsection{The prior, as coded}\label{sec:prior}

\emph{The prior has three layers: a prior on which predictors are in the model, a Gaussian prior on their
coefficients, and an inverse-gamma prior on the noise variance. Nothing in it is tuned by the user.}

\paragraph{Model space.} The method considers every model with at most $K$ predictors,
\begin{equation}\label{eq:modelspace}
  \Om:=\{S\subseteq\set{p}:\abs{S}\le K\},\qquad K:=\min\big(p,\lfloor n/2\rfloor\big).
\end{equation}

\paragraph{Model prior.} The prior probability of a model depends only on its size:
\begin{equation}\label{eq:modelprior}
  \varpi(S)\ \propto\ w(\abs{S}),\qquad w(k):=\frac{1}{2\binom{2p-1}{k}},\qquad S\in\Om.
\end{equation}
This is the beta-binomial prior $\BB(1,p)$ \citep{ScottBerger2010}, restricted to models of size at most $K$: each
predictor enters the model independently with the same probability $q$, and $q$ itself is drawn from $\BetaDist(1,p)$.
Its prior mean model size is about one. Adding one predictor to a model of size $k$ multiplies the prior weight by
$(k+1)/(2p-1-k)$, that is, it costs about $\log\{2p/(k+1)\}$ in log prior odds.

\begin{example}[The model prior with three predictors]\label{ex:prior}
Take $p=3$ and $n\ge6$, so that $K=3$ and no model is excluded. Then $w(0)=\tfrac12$, $w(1)=\tfrac1{10}$,
$w(2)=w(3)=\tfrac1{20}$. The empty model has prior probability $\tfrac12$; each one-predictor model has
$\tfrac1{10}$ (together $\tfrac3{10}$); each two-predictor model has $\tfrac1{20}$ (together $\tfrac3{20}$); and the full
model has $\tfrac1{20}$. The probabilities add up to one, so here the cap does not bind. Small models are favoured, and
each extra predictor must earn its place through the data.
\end{example}

\paragraph{Coefficients and noise variance.} Given a model $S$ and the noise variance $\sigma^2$, the coefficients of the
predictors in $S$ have the ridge-stabilised $g$-prior \citep{Zellner1986} with $g=n$, and the noise variance has an
inverse-gamma prior:
\begin{equation}\label{eq:slab}
  \beta_S\mid S,\sigma^2\ \sim\ \Normal\big(0,\;n\sigma^2(G_S+I)^{-1}\big),\qquad
  \sigma^2\ \sim\ \IG(a_0,b_0),\quad a_0=b_0=0.01.
\end{equation}
The working likelihood is $y\mid S,\beta_S,\sigma^2\sim\Normal_n(X_S\beta_S,\sigma^2I_n)$, on the standardised data. The
intercept is removed by centring and has no prior of its own. The prior and this likelihood are \emph{working}
assumptions of the method; the theorems are proved under the true model~\eqref{eq:model}.

\subsection{The posterior}\label{sec:posterior}

Integrating out $\beta_S$ and $\sigma^2$ gives the marginal likelihood of each model in closed form:
\begin{equation}\label{eq:marglik}
  m(S)\ \propto\ \Big(\frac{\det H_S}{\det(G_S+I)}\Big)^{-1/2}\,(2b_0+Q_S)^{-\nu},\qquad
  H_S:=(n+1)G_S+I,\quad Q_S:=y^\top y-n\,z_S^\top H_S^{-1}z_S,
\end{equation}
where $z_S:=X_S^\top y$ and $\nu:=a_0+n/2$. In words, $Q_S$ is a shrunken residual sum of squares of model $S$, and the
determinant ratio is an Occam factor that charges about $(n+1)^{1/2}$ for each predictor in a well-conditioned model.
The \emph{posterior over models} is
\begin{equation}\label{eq:posterior}
  \pis{S}=\frac{w(\abs{S})\,m(S)}{\sum_{V\in\Om}w(\abs{V})\,m(V)},\qquad S\in\Om,
\end{equation}
and $F(S):=\log w(\abs{S})+\log m(S)$ is the \emph{log score} of $S$: the log posterior, up to a constant. The method's
target is exactly this posterior. Because $y$ is standardised, the posterior over models does not change if $\yr$ is
shifted or rescaled.

\paragraph{Posterior summaries.} These are properties of the posterior, not of any algorithm:
\begin{itemize}[leftmargin=1.5em]
  \item the \emph{posterior inclusion probability} of predictor $j$, $\PIP_j:=\sum_{S\ni j}\pis{S}$;
  \item the \emph{highest posterior model} $\HPM:=\argmax_{S\in\Om}\pis{S}$;
  \item the \emph{median probability model} $\MPM:=\{j:\PIP_j>1/2\}$ \citep{BarbieriBerger2004};
  \item the \emph{posterior miss} $\miss:=1-\pis{\Sz}$, the posterior mass on wrong models (with $\pis{\Sz}:=0$ if
    $\Sz\notin\Om$).
\end{itemize}
Two elementary facts link them. For every data set, every inclusion probability is within $\miss$ of the truth,
$\max_j\abs{\PIP_j-\ind{j\in\Sz}}\le\miss$; and if $\miss<1/2$ then $\HPM=\MPM=\Sz$. So a bound on the posterior miss
controls all three summaries at once.

\subsection{What the method returns}\label{sec:method}

\emph{The method is a sequential Monte Carlo sampler over $\Om$ with a deterministic start, an internal calibration,
and a local search. Four of its objects appear in the theorems.}

\paragraph{The start law.} The method partitions the predictors into blocks $C_1,\dots,C_J$ of at most $16$ columns
each, by clustering on sample correlations; the theory uses only that the partition is a function of $X$ with blocks of
size at most $16$. Within each block it scores every subset exactly. The \emph{start law} is the posterior with all
interactions between blocks ignored:
\begin{equation}\label{eq:startlaw}
  \pi_0(S)\ \propto\ w(\abs{S})\prod_{b=1}^{J}\frac{m(S\cap C_b)}{m(\emptyset)},\qquad S\in\Om.
\end{equation}
It can be computed and sampled exactly. The \emph{start mode} $\tSz$ is a maximiser of $\pi_0$ over $\Om$.

\paragraph{Local search.} Starting at $\tSz$, the method repeatedly scans every \emph{neighbour} of the current model
(every deletion of one predictor, every addition of one predictor while the size is below $K$, and every swap of one
member for one non-member) and moves to the best neighbour if it strictly improves the log score $F$. It stops when a
full scan finds no strict improvement, or after at most $\ell_n+\abs{\tSz}$ moves, where
$\ell_n:=\lceil\log_2(n+p+1)\rceil$.

\paragraph{The reported model.} The method records every model it scores, in the sampler and in the local searches.
The \emph{reported model} $\Shat$ is the best-scoring recorded model. The method \emph{certifies} $\Shat$ (it sets its
local-optimality flag) when either the model space was small enough to be enumerated, so that $\Shat$ maximises $F$
over $\Om$, or $\Shat$ is the end point of a local search whose last full scan found no strict improvement, so that
$\Shat$ is a local maximum of $F$ for the neighbourhood above.

\paragraph{Estimated summaries.} The sampler runs $R=4$ independent \emph{islands}. Island $r$ produces estimated
inclusion probabilities $\PIPhat_{r,j}$ and an estimate $\Zhat_r$ of the normalising constant
$Z:=\sum_{V\in\Om}w(\abs{V})\,m(V)$ of the posterior (the \emph{evidence}). The method pools the islands with equal
weights,
\[
  \PIPhat_j:=\frac1R\sum_{r=1}^R\PIPhat_{r,j},\qquad \Zhat:=\frac1R\sum_{r=1}^R\Zhat_r,
\]
and reports the estimated median probability model $\MPMhat:=\{j:\PIPhat_j>1/2\}$. A run either returns these outputs
(it \emph{produces output}) or stops without them, for example when its computing allocation is exhausted.

\subsection{Standing assumptions}\label{sec:standing}

\emph{Two assumptions are used by every theorem. Three more are conventions about the computation, used only where a
theorem is about what the code does.}

\begin{assumption}{GL}{Gaussian linear model}
The data follow model~\eqref{eq:model}, with a fixed noise variance $\sigs^2>0$ and noise independent of the design.
\rulesout{heavy-tailed or heteroscedastic noise, and noise that depends on the predictors.}
\end{assumption}

\begin{assumption}{DIM}{Dimension}
$n\to\infty$ and $p=p_n$ satisfies $2\le p\le n^A$ for a fixed $A\ge1$. Many more predictors than observations
($p\gg n$) are allowed and are the main case.
\rulesout{a number of predictors growing faster than every polynomial in $n$.}
\end{assumption}

\begin{assumption}{MC}{The method as coded}
Statements about what the method returns assume exact arithmetic (scores and comparisons computed without rounding
error), the shipped configuration (in particular, the local search from the start mode is switched on with the move cap
above), and admissible runs (valid inputs and seed). Unless a statement says otherwise, it holds for every admissible
seed and every outcome of the method's internal calibration.
\rulesout{nothing statistical; it fixes what ``the method'' means.}
\end{assumption}

\begin{assumption}{DC}{Delivery conventions}
Used only by the clauses ``the run always produces output''. An idealised machine, the covered configuration, ideal
random numbers, and a declared memory and time allocation at least the method's polynomial envelopes in $(n,\praw)$.
The full list is in Appendix~\ref{app:conventions}. Under (DC), a run produces output on the bounded domain $\Dom$ of
$(n,\praw)$ on which the code's two fixed integer limits are not reached; it produces output for every $n$ under the
stronger convention (DC$^+$), which lifts those two limits.
\rulesout{nothing statistical; these are operational assumptions about the computer.}
\end{assumption}

\begin{assumption}{SC}{Sampler conventions}
Used only by the results on the printed estimates (Theorem~\ref{thm:outputs} and Remark~\ref{rem:unbiased}). Exact
arithmetic; ideal independent uniform random numbers, with the final run's random numbers independent of the data and of
the calibration; exact draws from the start law; Markov moves that realise their stated kernels. The full list is in
Appendix~\ref{app:conventions}.
\rulesout{nothing statistical; these idealise the random-number generator and the coded kernels.}
\end{assumption}

\subsection{Asymptotic conventions}\label{sec:asymptotics}

All results are asymptotic in $n\to\infty$ along a sequence of problems: $p=p_n$, the design, $\bstar$ and $\Sz$ may
change with $n$; the true model size $\sz$ is fixed unless a statement says otherwise. The reference rate is
$\rate:=\rn$.
For random variables $V_n$ and positive numbers $a_n$: $V_n=\OP(a_n)$ means that for every $\epsilon>0$ there is $C$ with
$\P(\abs{V_n}>Ca_n)\le\epsilon$ for all large $n$; $V_n=\oP(a_n)$ means $V_n/a_n\to0$ in probability. ``With probability
tending to one'' means $\P(\cdot)\to1$. For positive sequences, $a_n=O(b_n)$ means $a_n\le Cb_n$ for a constant $C$;
$a_n=o(b_n)$ means $a_n/b_n\to0$; $a_n\lesssim b_n$ is $a_n=O(b_n)$; $a_n\asymp b_n$ means $a_n\lesssim b_n$ and
$b_n\lesssim a_n$; $a_n\gg b_n$ means $b_n/a_n\to0$.
Constants written with Greek or Latin letters in an assumption ($\varphi_0$, $g_0$, $\bar\theta$, \dots) are fixed: they
do not depend on $n$. $C$ and $c$ denote positive constants that depend only on such fixed quantities; their value may
change from line to line.

% =====================================================================================================
\section{The posterior selects the right model}\label{sec:selection}

\emph{This section shows that the posterior puts almost all its mass on the true model, at the rate $\rn$, both for
fixed and for random Gaussian designs; that this rate cannot be improved in general; and that the same holds when the
fitted predictors omit some of the true ones.}

\subsection{Fixed designs}\label{sec:fixed}

We first state the assumptions on the design and the coefficients. Each comes with one line on what it rules out.

\begin{assumption}{SP}{Sparsity}
The true model size $\sz\ge1$ is fixed.
\rulesout{true models whose size grows with $n$ (see Theorem~\ref{thm:sharp} for what is possible then).}
\end{assumption}

\begin{assumption}{BC}{Bounded coefficients}
$\max_j\abs{\theta^*_j}\le\bar\theta$ for a fixed $\bar\theta<\infty$.
\rulesout{unbounded signal energy, under which the posterior of the noise variance could drift off scale.}
\end{assumption}

\begin{assumption}{SE}{Sparse eigenvalue floor}
$\lmin(\Rhat_S)\ge\varphi_0^2$ for every $S\subseteq\set{p}$ with $\abs{S}\le2(A+1)\sz+2$, for a fixed
$\varphi_0\in(0,1]$.
\rulesout{strong collinearity within small sets of predictors, which would let a wrong small model mimic the true one.}
\end{assumption}

\begin{assumption}{NC}{No near-collinearity among $O(\log n)$ columns}
$\lmin(X_S^\top X_S)\ge g_0$ for every $S$ with $\abs{S}\le3(A+1)(\sz+1)\log n$, for a fixed $g_0>0$.
\rulesout{exact or near-exact linear dependence among $O(\log n)$ columns, duplicated columns included. On the
correlation scale the floor is $g_0/(n-1)$, so it vanishes as $n$ grows.}
\end{assumption}

\begin{assumption}{BM}{Beta-min}
Either (a) $\tmin\ge C_\beta\,\sigs\rn$ with a fixed $C_\beta>0$ such that $C_\beta^2\varphi_0^2>c_0(\sz,A)$, where
$c_0(\sz,A)$ is an explicit constant depending only on $(\sz,A)$; or (b) $\tmin\sqrt{n/\log n}\to\infty$.
\rulesout{true signals too small to be told apart from noise at the scale $\rn$.}
\end{assumption}

\inwords{If the true model is small, the coefficients are bounded, small groups of predictors are not collinear, and
every true signal is at least a constant times $\rn$, then the posterior puts mass at least $1-C\rn$ on the true model,
on average over the noise, uniformly over all such designs. As a consequence, the highest posterior model and the
median probability model both equal the true model with probability tending to one, and every inclusion probability is
within $\OP(\rn)$ of $0$ or $1$.}

\begin{theorem}[Posterior concentration, fixed design]\label{thm:post-fixed}
Let the design be fixed and assume \as{GL}, \as{DIM}, \as{SP}, \as{BC}, \as{SE}, \as{NC} and \as{BM}. There are $n_0$
and $C$, depending only on $(\sz,A,\sigs^2,\bar\theta,\varphi_0,g_0,C_\beta)$, such that for every $n\ge n_0$ and every
design and $\tstar$ that satisfy these assumptions at that $n$,
\begin{equation}\label{eq:post-exp}
  \E\big[1-\pis{\Sz}\big]\ \le\ C\sqrt{\frac{\log n}{n}} .
\end{equation}
Consequently, along every sequence that satisfies the assumptions,
\begin{gather}
  \P\big(\HPM=\MPM=\Sz\big)\ \ge\ 1-4C\sqrt{\frac{\log n}{n}}\ \to1,\label{eq:post-cons}\\
  \max_{j\le p}\big|\PIP_j-\ind{j\in\Sz}\big|\ \le\ 1-\pis{\Sz}\ =\ \OP\Big(\sqrt{\frac{\log n}{n}}\Big).\nonumber
\end{gather}
The expectation is over the noise, with the design fixed. Under the second form of \as{BM}, $n_0$ also depends on how
fast $\tmin\sqrt{n/\log n}$ diverges.
\end{theorem}

\status{The in-probability conclusions~\eqref{eq:post-cons}: \textbf{proved and independently checked}. The expectation
bound~\eqref{eq:post-exp}: \textbf{proved and checked, conditional on two supporting lemmas} (a tail bound on the
posterior model size and a lower bound on the posterior miss) \textbf{that are proved but not yet independently
checked}. \bridgepending{} Verbatim registry strings: Appendix~\ref{app:status}, row 1.1.}

The next result drops the eigenvalue floor away from the true model and the beta-min condition, and replaces both by a
single identifiability margin.

\begin{assumption}{ID}{Floor on the truth and an identifiability margin}
$\lmin(\Rhat_{\Sz})\ge\varphi_0^2$, and every model $S$ with $\abs{S}\le K$, $\Sz\not\subseteq S$ and at most
$j_0:=\lceil(2A+1)\sz\rceil+1$ predictors outside $\Sz$ leaves residual signal
$\norm{(I-P_S)X\tstar}^2\ge c\,\sigs^2\log n$, for a fixed $c>c_1(\sz,A)$, where $c_1(\sz,A)$ is an explicit constant
depending only on $(\sz,A)$.
\rulesout{small wrong models that explain almost all of the mean signal; it does not restrict correlations among the
other predictors.}
\end{assumption}

\inwords{The same conclusions hold if, instead of a floor on every small set of predictors and a lower bound on each true
coefficient, we only ask that the true predictors are not collinear among themselves and that every small model that
misses a true predictor leaves a residual signal of order $\log n$. The vanishing floor \as{NC} is kept.}

\begin{theorem}[Posterior concentration under an identifiability margin]\label{thm:post-margin}
Let the design be fixed and assume \as{GL}, \as{DIM}, \as{SP}, \as{BC}, \as{NC} and \as{ID}. Then the conclusions of
Theorem~\ref{thm:post-fixed} hold, with constants depending on $(\sz,A,\sigs^2,\bar\theta,\varphi_0,g_0,c)$.
\end{theorem}

\status{\textbf{Proved and checked, conditional on the same two supporting lemmas} as the expectation bound of
Theorem~\ref{thm:post-fixed}. \bridgepending{} Verbatim: Appendix~\ref{app:status}, row 1.1.}

\begin{remark}[On the assumptions]\label{rem:post-fixed}
(i) \emph{Sensitivity and specificity.} The bound $\max_j\abs{\PIP_j-\ind{j\in\Sz}}\le1-\pis{\Sz}$ holds for every data
set. With~\eqref{eq:post-exp} it gives both sensitivity and specificity of all inclusion probabilities at rate $\rn$,
uniformly in $j$.
(ii) \emph{Without beta-min.} Under the other assumptions of Theorem~\ref{thm:post-fixed}, the posterior mass of models
that strictly contain $\Sz$ is $\OP(\rn)$, and every true predictor with $\theta_j^{*2}\ge c\,\sigs^2\log n/(n\varphi_0^2)$,
$c>c_0(\sz,A)$, has $1-\PIP_j=\OP(\rn)$. Specificity can fail without beta-min (numerical counterexample).
(iii) \emph{Constants.} $c_0(\sz,A)$ and $c_1(\sz,A)$ are explicit but not optimised. Under \as{SE} and \as{BM}(a) the
margin \as{ID} holds with $c=C_\beta^2\varphi_0^2(1-1/n)$, so Theorem~\ref{thm:post-margin} contains
Theorem~\ref{thm:post-fixed} when $C_\beta^2\varphi_0^2>c_1$; for $c_0<C_\beta^2\varphi_0^2\le c_1$ the two theorems
coexist.
(iv) \emph{Coefficients.} \as{BC} is used only through the energy bound $\norm{X\tstar}^2\le n\sz^2\bar\theta^2$; any
bound $\norm{X\tstar}^2\le nM$ with a fixed $M$ suffices.
(v) \emph{Large models.} Models of size at least $3(A+1)(\sz+1)\log n$ carry posterior mass $\OP(n^{-1/2})$ with no
eigenvalue condition.
(vi) \emph{A weaker form of \as{NC}.} \as{NC} can be weakened to a condition on the Occam factors of over-fitted models;
a readable sufficient form is determinant growth,
$\det H_{\Sz\cup J}\det(G_{\Sz}+I)/\{\det H_{\Sz}\det(G_{\Sz\cup J}+I)\}\ge\{c_D(n+1)\}^{\abs{J}}$ for
$1\le\abs{J}\le3(A+1)(\sz+1)\log n$.
(vii) \emph{Necessity.} The floor on the true model in \as{ID} is sufficient and not known to be necessary.
Theorem~\ref{thm:post-margin} keeps the vanishing floor \as{NC}, which still excludes exact linear dependence among
$O(\log n)$ columns, duplicated columns included.
\end{remark}

\begin{remark}[What the posterior looks like at moderate $n$]\label{rem:moderate}
The theorems describe the limit $n\to\infty$. The miss $1-\pis{\Sz}$ is carried mainly by the $p-\sz$ models
$\Sz\cup\{j\}$ that add one null predictor to the truth. Each is individually light: its posterior odds against $\Sz$
are of order $n^{-1/2}/p$, times a $\chi^2_1$ factor. But they are numerous, and together they can weigh as much as
$\Sz$ itself. So at moderate $n$ and large $p$ the posterior is a mixture over $\Sz$ and many one-predictor supersets
of it.
\end{remark}

\begin{remark}[Comparison with \citet{SongLiang2015}]\label{rem:SL-fixed}
Theorems~\ref{thm:post-fixed} and~\ref{thm:post-margin} give full-model selection consistency for the shipped prior,
with an explicit, design-uniform expectation rate and exact HPM/MPM; no tuning (Song--Liang's Theorem~3 gives sure
screening only: sensitivity $\to1$, specificity not controlled; their MAP consistency is in their supplement).
Theorem~\ref{thm:post-margin} replaces the bounded-size floor \as{SE} off the truth and beta-min by a residual-signal
identifiability margin, while retaining the vanishing floor \as{NC} (in the style of \citealp{NarisettyHe2014}; compare
Song--Liang's population-level identifiability Condition~3 for their Theorem~3). Narrower in dimension ($n^A$ against
Song--Liang's $P_n<e^{Cn^\alpha}$, $\alpha<1$: their Corollary~1 via their~(10), inherited by their Theorem~3).
\end{remark}

\subsection{Random Gaussian designs}\label{sec:random}

\emph{When the rows of the design are drawn from a Gaussian law, the sample conditions of Section~\ref{sec:fixed} hold
with probability tending to one under a condition on the population correlation matrix alone. The conclusions are the
same.}

\begin{assumption}{GR}{Gaussian rows}
The rows of $\Xr$ are i.i.d.\ $\Normal_p(m_n,\Sig)$, with $\Sig$ a correlation matrix.
\rulesout{dependent rows and non-Gaussian predictors (sub-Gaussian rows are expected to work but are not claimed).}
\end{assumption}

\begin{assumption}{FS}{Fixed signals}
$b_{\mathrm{lo}}\le\abs{\beta^*_j}\le b_{\mathrm{hi}}$ for every $j\in\Sz$, with fixed $0<b_{\mathrm{lo}}\le b_{\mathrm{hi}}<\infty$.
\rulesout{true coefficients that shrink to zero or blow up with $n$.}
\end{assumption}

\begin{assumption}{PE}{Population sparse eigenvalue floor}
$\lmin(\Sigma_{n,S})\ge\varphi^2$ for every $S$ with $\abs{S}\le5(A+1)(\sz+1)\log n$, for a fixed $\varphi\in(0,1]$ (for
instance, $\lmin(\Sig)\ge\varphi^2$).
\rulesout{population collinearity within sets of $O(\log n)$ predictors. No upper bound on the eigenvalues is needed.}
\end{assumption}

\inwords{If the rows are Gaussian, the true coefficients are fixed and nonzero, and no set of $O(\log n)$ predictors is
collinear in the population, then the posterior concentrates on the true model at the rate $\rn$, exactly as for fixed
designs, now under the joint law of the design and the noise.}

\begin{theorem}[Posterior concentration, random Gaussian design]\label{thm:post-random}
Assume \as{GL}, \as{DIM}, \as{SP}, \as{GR}, \as{FS} and \as{PE}. Then, under the joint law of the design and the noise,
\begin{gather*}
  1-\pis{\Sz}=\OP\Big(\sqrt{\frac{\log n}{n}}\Big),\qquad \P\big(\HPM=\MPM=\Sz\big)\to1,\\
  \max_j\big|\PIP_j-\ind{j\in\Sz}\big|=\OP\Big(\sqrt{\frac{\log n}{n}}\Big).
\end{gather*}
\end{theorem}

\status{\textbf{Proved and independently checked} (the theorem it rests on and the transfer from population to sample
conditions are both checked). \bridgepending{} Verbatim: Appendix~\ref{app:status}, row 1.2.}

\begin{remark}\label{rem:post-random}
(i) No upper bound on the eigenvalues of $\Sig$ and no beta-min condition beyond \as{FS} are needed, and $p\gg n$ is
allowed. The underlying result allows the floor in \as{PE} to vanish slowly, as long as it is much larger than
$(\log n)^{3/2}/\sqrt n$; the statement above uses a constant floor for simplicity.
(ii) Gaussian rows are used only to show that sample correlations of small sets of predictors are close to their
population values. Without the lower bound $b_{\mathrm{lo}}$ in \as{FS}, the statements of
Remark~\ref{rem:post-fixed}(ii) still hold.
(iii) This is the random-design setting of \citet{SongLiang2015}, here with $p\gg n$.
\end{remark}

\subsection{The rate cannot be improved}\label{sec:sharp}

\emph{The rate $\rn$ of the two previous theorems is the right one. On every design the posterior miss is at least of
order $n^{-1/2}$, and on designs whose null predictors are orthogonal or independent Gaussian it is at least of order
$\rn$, so the extra $\sqrt{\log n}$ is really there.}

The lower bounds use the following conditions.

\begin{assumption}{EN}{Energy bounds}
$\norm{X\tstar}^2\le nM$ and $\norm{\tstar}^2\le nM'$ for fixed $M,M'$.
\rulesout{signal energy growing faster than $n$.}
\end{assumption}

\begin{assumption}{NO}{Orthogonal null columns}
The columns of $X$ outside $\Sz$ are mutually orthogonal and orthogonal to the columns in $\Sz$.
\end{assumption}

\begin{assumption}{NG}{Gaussian null columns}
The raw columns outside $\Sz$ have i.i.d.\ $\Normal(0,1)$ entries, independent of $\Xr_{\Sz}$ and of $\varepsilon$, and
$p\le n^A$.
\end{assumption}

Assumptions \as{NO} and \as{NG} are not restrictions we need for consistency; they describe the two \emph{witness}
designs on which we show that the rate is attained.

\inwords{(a) Whatever the design, the posterior cannot put more than $1-c/\sqrt n$ of its mass on any one small model,
because the prior and the Occam factor keep every one-predictor extension of it alive. (b) On designs whose null
predictors are orthogonal, or independent Gaussian, the posterior miss is at least a constant times
$(\sz+1)\sqrt{\log n/n}$ on average, for every number of predictors. So the rate $\rn$ of
Theorem~\ref{thm:post-fixed} cannot be improved under its assumptions.}

\begin{theorem}[The rate cannot be improved]\label{thm:sharp}
\textup{(a)} \emph{(Every design.)} For every $n\ge3$, every standardised design with $p\ge2$, every response and every
model $S$ with $\abs{S}\le\min(K-1,p/2)$,
\begin{equation}\label{eq:sharp-a}
  \pis{S}\ \le\ \Big(1+\frac{\abs{S}+1}{3\sqrt n}\Big)^{-1}.
\end{equation}
\textup{(b)} \emph{(Orthogonal or Gaussian null columns.)} Assume \as{GL}, that $\sz$ is fixed, \as{EN}, and either
\as{NO} or \as{NG}. Then, for every $p$ with $\sz<p$, bounded or growing,
\begin{equation}\label{eq:sharp-b}
  \E\big[1-\pis{\Sz}\big]\ \ge\ (1-o(1))\,(\sz+1)\,\frac{2(p-\sz)}{2p-1-\sz}\,
  \sqrt{\frac{\log\!\big(p\sqrt n/(\sz+1)\big)}{\pi n}},
\end{equation}
and, if moreover $p\to\infty$,
\begin{equation}\label{eq:sharp-c}
  1-\pis{\Sz}\ \ge\ (1-\oP(1))\,(\sz+1)\sqrt{\frac{\log p}{\pi n}}.
\end{equation}
\end{theorem}

In~\eqref{eq:sharp-b} and~\eqref{eq:sharp-c}, $\pi$ is the number $3.14\ldots$, not the posterior. Display~\eqref{eq:sharp-a}
says that no model, the true one included, can carry posterior mass closer to one than about $(\abs{S}+1)/(3\sqrt n)$.
Display~\eqref{eq:sharp-b} says that, on the witness designs, the expected posterior miss is at least
$c\cdot(\sz+1)\cdot\rn$, where $c>0$ depends only on $\sz$, and does not depend on $\sz$ when $\sz\le p/2$.

\status{\textbf{Proved and independently checked}. \bridgepending{} Verbatim: Appendix~\ref{app:status}, row 1.3.}

\begin{remark}[What is attained where]\label{rem:sharp}
(i) By~(a), for true models with $\sz\le\min(K-1,p/2)$, no design gives a rate faster than $n^{-1/2}$;
$\pis{\Sz}\to1$ requires $\sz=o(\sqrt n)$; and the rate $\rn$ requires $\sz=O(\sqrt{\log n})$. (Without the size
restriction, (a) is false: for example $S=\set{p}$ with $p=2$.)
(ii) By~(b), on orthogonal or Gaussian null columns the factor $\sqrt{\log n}$ is present \emph{in expectation for
every $p$}, bounded $p$ included (with $p=n^a$ the constant is $\sqrt{(a+\frac12)/\pi}$), and \emph{in probability
whenever $\log p\asymp\log n$}. The extra $\sqrt{\log n}$ in expectation at bounded $p$ comes from rare large gains of a
single extra predictor, not from the number of competitors.
(iii) \emph{Exactness.} The rate $\rn$ of Theorem~\ref{thm:post-fixed} is exact in two named cases. In expectation: on
fully orthogonal designs (all columns mutually orthogonal, any $p$ with $\sz<p\le n-1$, bounded or growing) with fixed
nonzero coefficients; these designs satisfy the assumptions of Theorem~\ref{thm:post-fixed} and of~\eqref{eq:sharp-b},
so $\E[1-\pis{\Sz}]$ is of exact order $\rn$, with constants depending only on $\sz$, the coefficients and the noise
variance (conditional on Theorem~\ref{thm:post-fixed}). In probability: on fully orthogonal designs with
$p=\min(n-1,\lfloor n^a\rfloor)$ and on raw designs with i.i.d.\ $\Normal(0,1)$ entries and $p=\lfloor n^a\rfloor$, both
with fixed nonzero coefficients (conditional on Theorems~\ref{thm:post-fixed} and~\ref{thm:post-random}).
(iv) \emph{Growing $\sz$.} Keep \as{EN}, let all columns outside $\Sz$ be null columns of type \as{NO}, or of type
\as{NG} with the columns in $\Sz$ and the coefficients fixed (nonrandom), and let $\sz\le p/2$ grow with $n$. Then a
posterior mean error $\E[1-\pis{\Sz}]$ of order $\rn$ forces $\sz$ to stay bounded, for every $p$.
(v) \emph{Critical scale} (orthogonal null columns only). Keep \as{EN} and \as{NO}, let $p\to\infty$ and let $\sz$ grow
with $\sz/p\to0$ and $(\sz+1)\sqrt{\log p/(\pi n)}\to\lambda\in(0,\infty)$. Then $\pis{\Sz}\le e^{-\lambda}+\oP(1)$.
(vi) \emph{What is open.} On a general design only the order $n^{-1/2}$ of part~(a) is proved as a lower bound; a
matching lower bound $\E[1-\pis{\Sz}]\gtrsim\rn$ under a floor on $\Sz\cup\{j\}$ is open. When $\log p=o(\log n)$,
nothing is claimed in probability: \eqref{eq:sharp-c} gives only the order $\sqrt{\log p/n}$, and whether the
$\sqrt{\log n}$ is present in probability there is open. On raw Gaussian designs the matching upper bound in expectation
is open (Theorem~\ref{thm:post-random} is an in-probability statement). Whether the constant in~\eqref{eq:sharp-b} is
attained is not proved.
\end{remark}

\subsection{Omitted variables}\label{sec:omitted}

\emph{In practice some true predictors may be missing from the data. Here the method is run on a set of candidate
predictors only. If the response, given the candidates, is still Gaussian and linear in them, the posterior
concentrates on the support of the best linear predictor built from the candidates, at the same rate $\rn$.}

In this subsection the method is run on the response $\yr$ and a set $\mathcal C$ of candidate predictors only, with
rows $x_{\mathcal C,1},\dots,x_{\mathcal C,n}$; $p$ is the number of candidates, $\Sigma$ the covariance matrix and $R$
the correlation matrix of a candidate row.

\begin{assumption}{CM}{Candidate model}
The candidate rows are i.i.d.\ $\Normal_p(m,\Sigma)$, and, given the candidate design, the responses are independent
with
\[
  \yr_i\mid x_{\mathcal C}\ \sim\ \Normal\big(\alpha'+x_{\mathcal C,i}^\top\beta^{\mathcal C},\ \sigma^2_{\mathcal C}\big),
  \qquad \sigma^2_{\mathcal C}>0 .
\]
We write $\Sz^{\mathcal C}:=\supp\beta^{\mathcal C}$ and $s:=\abs{\Sz^{\mathcal C}}$.
\rulesout{a conditional law of the response, given the candidates, that is not Gaussian with a mean linear in the
candidates and a constant variance.}
\end{assumption}

In words, $\beta^{\mathcal C}$ are the coefficients of the best linear predictor of the response from the candidates,
and $\sigma^2_{\mathcal C}$ is the variance left over, which includes the contribution of the omitted predictors.

\begin{assumption}{CF}{Candidate floor}
$\lmin(R_{TT})\ge c$ for every $T\subseteq\mathcal C$ with $\abs{T}\le3(A+1)(s+1)\log n$, for a fixed $c>0$.
\rulesout{population collinearity among $O(\log n)$ candidates.}
\end{assumption}

\begin{assumption}{CS}{Signal relative to the residual variance}
$\beta^{\mathcal C\top}\Sigma\beta^{\mathcal C}\le M_0\,\sigma^2_{\mathcal C}$, and
\[
  \frac{n\min_{j\in\Sz^{\mathcal C}}\Sigma_{jj}(\beta^{\mathcal C}_j)^2}{\sigma^2_{\mathcal C}\log n}\to\infty .
\]
\rulesout{candidate signal that is huge compared with the residual variance, and candidate coefficients too small to be
detected at the scale $\rn$.}
\end{assumption}

\inwords{If the method sees only some of the predictors, and the response given those predictors is Gaussian and linear
in them, the posterior picks the predictors that the best linear predictor uses, with error of order $\rn$. The
residual variance may be fixed, may vanish, or may grow with $n$.}

\begin{theorem}[Omitted variables]\label{thm:omitted}
Assume \as{CM}, \as{CF}, \as{CS}, \as{DIM}, and that $s\ge1$ is fixed. No condition is placed on how $m$, $\Sigma$,
$\beta^{\mathcal C}$ or $\sigma^2_{\mathcal C}$ vary with $n$. Then
\[
  \E\big[1-\pi(\Sz^{\mathcal C}\mid y)\big]=O\Big(\sqrt{\frac{\log n}{n}}\Big),\qquad
  \P\big(\HPM\neq\Sz^{\mathcal C}\big)=O\Big(\sqrt{\frac{\log n}{n}}\Big).
\]
\end{theorem}

\status{\textbf{Proved and checked, conditional on} the expectation bound of Theorem~\ref{thm:post-fixed} (hence on its
two supporting lemmas, which are not yet independently checked) and on two lemmas on Gaussian designs that are verified
and passed a council review. \bridgepending{} Verbatim: Appendix~\ref{app:status}, row 1.4.}

\begin{example}[Omitted Gaussian predictors]\label{ex:omitted}
Suppose the rows $(x_{\mathcal C,i},x_{\mathcal M,i})$ of candidates and omitted predictors are i.i.d.\ Gaussian with
$\Sigma_{\mathcal C\mathcal C}$ invertible, and
$\yr_i=\alpha+x_{\mathcal C,i}^\top\beta_{\mathcal C}+x_{\mathcal M,i}^\top\beta_{\mathcal M}+\varepsilon_i$ with
$\varepsilon\sim\Normal_n(0,\sigma^2I)$ independent of the rows. Then the full model is misspecified for the
candidate-only fit, and \as{CM} holds with
\[
  \beta^{\mathcal C}=\beta_{\mathcal C}+\Sigma_{\mathcal C\mathcal C}^{-1}\Sigma_{\mathcal C\mathcal M}\beta_{\mathcal M},
  \qquad
  \sigma^2_{\mathcal C}=\sigma^2+\beta_{\mathcal M}^\top\big(\Sigma_{\mathcal M\mathcal M}
  -\Sigma_{\mathcal M\mathcal C}\Sigma_{\mathcal C\mathcal C}^{-1}\Sigma_{\mathcal C\mathcal M}\big)\beta_{\mathcal M}.
\]
The first formula says that each candidate absorbs the part of the omitted signal it can predict; the second that the
rest of the omitted signal is added to the noise. If $\Sigma_{\mathcal C\mathcal C}$ is singular, the same holds with its
Moore--Penrose inverse. The same holds when the omitted predictors are a linear function of the candidates plus Gaussian
noise whose rows are i.i.d.\ and independent of the candidates.
\end{example}

\begin{remark}[On Theorem~\ref{thm:omitted}]\label{rem:omitted}
(i) \emph{What the posterior targets.} The law in \as{CM} is the conditional law of $\yr$ given the candidates, hence
the Kullback--Leibler projection of the truth onto the candidate linear models: the posterior concentrates on the
support of the best candidate predictor. Under the same assumptions the posterior's fitted mean and variance also
contract around this law, with its mean centred over the sample (the fitted mean carries no intercept), in the
in-sample average Hellinger and Kullback--Leibler distances, at rate $M_n/\sqrt n$ in Hellinger and $M_n/n$ in
Kullback--Leibler, for any $M_n\to\infty$. For a fixed design the contraction needs only that the posterior mass of
$\Sz^{\mathcal C}$ tends to one in probability (with the energy bound and a floor on $\Sz^{\mathcal C}$); it holds
uniformly over any class of fixed-design configurations with common constants and bounded support size on which this
convergence is uniform. A fixed-design version of the theorem, uniform over the residual variance, also holds.
(ii) \emph{A weaker sufficient condition.} The proof uses only that, given the candidates, the coordinates of the
residual vector $\yr-\alpha'\one-x_{\mathcal C}\beta^{\mathcal C}$ in an orthonormal basis of the vectors orthogonal to
$\one$ are i.i.d.\ $\Normal(0,\sigma^2_{\mathcal C})$; nothing is required of the residual's sample mean. The Gaussian law
of the candidate rows is used only to transfer the population floor \as{CF} and \as{CS} to the sample.
(iii) \emph{Why relative to $\sigma^2_{\mathcal C}$.} Everything is measured relative to $\sigma^2_{\mathcal C}$ because
the method is invariant to shifts and rescaling of $y$. The energy bound in \as{CS} cannot simply be dropped (numerical
evidence); the beta-min order is that of Theorem~\ref{thm:post-fixed}.
(iv) \emph{Relation to Theorem~\ref{thm:post-random}.} Theorem~\ref{thm:post-random} does not give
Theorem~\ref{thm:omitted}: it fixes the noise variance and has no omitted predictors. Whether the hypotheses of
Theorem~\ref{thm:omitted} cover the setting of Theorem~\ref{thm:post-random} without omitted predictors is open; until
then the two theorems rest on different arguments, each with its own conditions.
\end{remark}

\status{Remark~\ref{rem:omitted}(i): \textbf{proved and checked}, conditional on the posterior contraction lemma for
fixed designs (itself proved and checked) and, for the random-design clause, on the same inputs as
Theorem~\ref{thm:omitted}. Verbatim: Appendix~\ref{app:status}, row 1.4.}

\begin{remark}[Comparison with \citet{SongLiang2015}]\label{rem:SL-omitted}
Song--Liang (2015, Theorem~4, eqs.~(17)--(18); Corollary~2) fit a fixed candidate subset that misses true predictors, a
model they call misspecified, with candidates having finite second moments only (their Condition~4): the posterior is
Hellinger-consistent at the Kullback--Leibler projection, for each fixed $\varepsilon$, almost surely and without a rate,
and, if the candidates have no exact collinearity, the posterior probability that a sampled model contains the
projection's support tends to one in probability (Corollary~2). Here, assuming instead that, given the candidates, the
response is Gaussian with a mean linear in them and a constant variance (as when the omitted predictors are jointly
Gaussian with the candidates), we prove exact selection of the best candidate linear predictor's support (Song--Liang's
$\tilde s$ for centred rows) at rate $\rn$, with up to $n^A$ candidates and no condition on how $\Sigma$, $\beta$,
$\sigma^2$ vary with $n$; the coded method is exactly invariant to shifts and rescaling of $y$.
\end{remark}

% =====================================================================================================
\section{The posterior is right about the distribution}\label{sec:distribution}

\emph{Selecting the right model is one goal; predicting well is another. This section shows that the posterior also
recovers the conditional law of the response given the predictors, at the rate $\rn$, in the distance used by
\citet{SongLiang2015}, with no beta-min condition and no fixed lower bound on eigenvalues. It then shows that the same
holds for the part of the posterior that a completed run of the sampler returns.}

\subsection{Setting}\label{sec:dist-setting}

In this section the data are $n$ independent copies $(x_i,y_i)$, $i=1,\dots,n$, of a pair $(x,y)$: $x\in\R^{\praw}$ is the
vector of all observed predictors (a row of $\Xr$) and $y\in\R$ the response (an entry of $\yr$). We write $D_n$ for the
data set. The number $\praw$ may grow with $n$ and may be much larger than $n$, and the law of $(x,y)$ may change with
$n$. As everywhere, the method runs on the $p\le\praw$ predictors that are not constant in the sample, with the prior of
Section~\ref{sec:prior}. We write $\Pi(\cdot\mid D_n)$ for the resulting exact posterior of $(S,\beta,\sigma^2)$, mapped
back to the original units: $S$ is the set of included predictors, $\beta\in\R^{\praw}$ the coefficient vector (zero
outside $S$ and on removed predictors), and $\sigma^2$ the noise variance. (If fewer than two predictors remain, or the
response is constant in the sample, the posterior is not defined; under the assumptions below such data sets have
probability smaller than the exceptional probabilities in the theorems and are counted in them, and $p=\praw$ outside
them.)

\paragraph{What is estimated.} A posterior draw $(S,\beta,\sigma^2)$ defines a conditional law of $y$ given $x$,
\[
  \Normal\big(\bar y+(x-\bar x)^\top\beta,\ \sigma^2\big),
\]
with the sample-mean intercept ($\bar y$ and $\bar x$ are sample means). The law of $x$ is \emph{not} estimated. Let
$f$ be the joint density of $(x,y)$ in which $x$ has its true law and $y$ given $x$ has the law above, and let $f^*$ be
the true joint density. For two laws, $H^2(P,Q):=1-\int\sqrt{dP\,dQ}$ is the squared Hellinger distance, and for two
densities $d_H(f,g):=(1-\int\sqrt{fg})^{1/2}$. Because $f$ and $f^*$ share the law of $x$,
\begin{equation}\label{eq:dH}
  d_H^2(f,f^*)=\E_x\Big[H^2\Big(\Normal\big(\bar y+(x-\bar x)^\top\beta,\sigma^2\big),\
  \Normal\big(\alpha^*+x^\top\bstar,\sigs^2\big)\Big)\Big],
\end{equation}
the squared Hellinger distance between the fitted and the true conditional laws of $y$ given $x$, averaged over a new
$x$ independent of the data. This is the distance of \citet[Theorems~1--2]{SongLiang2015}, who use
$\int(\sqrt f-\sqrt{f^*})^2=2d_H^2$.

\paragraph{Further notation.} $\P^*$ is the law of the data. For the predictors, $\sd$, $\Var$ and $\Corr$ are population
standard deviations, variances and correlation matrices; $z_j:=(x_j-\E x_j)/\sd(x_j)$ and
$z:=(z_1,\dots,z_{\praw})$; $x_T$ is the sub-vector of $x$ indexed by a set $T$.

\paragraph{Constants, level and window.} $s\ge0$ is an integer and $A\ge1$, $L\ge1$, $c_\psi>0$, $\bar b>0$ are
constants; none depends on $n$. The reader fixes a \emph{level} $M\ge1$: a larger $M$ gives a larger ball and a smaller
exceptional probability. The \emph{window} is
\[
  q_n:=\min\big\{\praw,\ \lceil C_qM\log n\rceil\big\},\qquad C_q:=2\big(s(A+2)+4\big)/\log2,
\]
the largest number of predictors that \as{SG} and \as{NE} below constrain jointly; it is of order $\log n$ for a fixed
level.

\subsection{Assumptions and results}\label{sec:dist-results}

\begin{assumption}{IID}{Gaussian linear model for i.i.d.\ pairs}
$y=\alpha^*+x^\top\bstar+\varepsilon$ with $\varepsilon\sim\Normal(0,\sigs^2)$, $\sigs^2>0$, independent of $x$. The
numbers $\alpha^*$, $\bstar$, $\sigs$ and the law of $x$ may depend on $n$.
\rulesout{non-Gaussian or heteroscedastic noise, and dependence between observations.}
\end{assumption}

\begin{assumption}{SD}{Sparsity and dimension}
At most $s$ coordinates of $\bstar$ are nonzero, and $2\le\praw\le n^A$.
\rulesout{dense coefficient vectors, and more than polynomially many predictors.}
\end{assumption}

\begin{assumption}{SG}{Sub-Gaussian predictors on sparse directions}
Every predictor has positive variance, and for every $u\in\R^{\praw}$ with at most $q_n$ nonzero entries and every
$\lambda\in\R$, $\E\exp(\lambda u^\top z)\le\exp\big(L^2\lambda^2\Var(u^\top z)/2\big)$.
\rulesout{heavy-tailed predictors, along sparse directions.}
\end{assumption}

\begin{assumption}{NE}{No near-exact collinearity}
$\lmin(\Corr(x_T))\ge c_\psi/n$ for every set $T$ of at most $q_n$ predictors.
\rulesout{groups of $O(\log n)$ predictors that are collinear up to $1-c_\psi/n$. Predictors may otherwise be very
strongly correlated.}
\end{assumption}

\begin{assumption}{BS}{Bounded signal}
$\abs{\beta^*_j}\sd(x_j)\le\bar b\,\sigs$ for every $j$.
\rulesout{a single predictor whose signal is unboundedly large compared with the noise, and huge coefficients that
cancel.}
\end{assumption}

No lower bound on the nonzero coefficients (no beta-min condition) and no fixed lower bound on eigenvalues is assumed.

\inwords{Under these assumptions, except on a set of data sets of probability at most $n^{-cM}$, the posterior gives
probability at most $n^{-cM}$ to draws whose conditional law of $y$ given $x$ is farther than $\sqrt{M\log n/n}$ from the
true one, in Hellinger distance averaged over $x$. That is, the posterior contracts around the true conditional law at
the rate $\rn$.}

\begin{theorem}[Distribution-level consistency]\label{thm:dist}
There are constants $c>0$, $C_1$, $M_0$ and $n_0$, depending only on $(s,A,L,c_\psi,\bar b)$, such that the following
holds. Fix a level $M\ge M_0$ and assume \as{IID}, \as{SD}, \as{SG}, \as{NE} and \as{BS}, with the window $q_n$ of this
$M$. Then for every $n\ge n_0$ with $\sqrt n/\log n\ge M$,
\begin{equation}\label{eq:dist-i}
  \P^*\Big[\,\Pi\Big(d_H(f,f^*)>\sqrt{M\log n/n}\ \Big|\ D_n\Big)>n^{-cM}\Big]\ \le\ n^{-cM}.
\end{equation}
\end{theorem}

\status{\textbf{Proved and independently checked}. Verbatim: Appendix~\ref{app:status}, row 2.1.}

\begin{remark}[Level-$M$ form, model size, and the case without \as{NE}]\label{rem:dist}
(i) \emph{Level-$M$ form.} The same constants serve every level $M$ with $M_0\le M\le\sqrt n/\log n$, which may depend
on $n$: whenever $n\ge n_0$ and the assumptions hold with the window $q_n$ of this $M$, \eqref{eq:dist-i} holds, and
also
\begin{equation}\label{eq:dist-ii}
  \P^*\Big[\,\Pi\big(\abs{S}>s+C_1M\ \big|\ D_n\big)>n^{-cM}\Big]\ \le\ n^{-cM}.
\end{equation}
The two exceptional sets of data sets can be taken to be the same set, and the bounds are uniform over all laws of the
data that satisfy the assumptions with the given constants. Display~\eqref{eq:dist-ii} bounds the number of selected
predictors; it does not say that they are the true ones.
(ii) \emph{Without \as{NE}.} If \as{NE} is dropped and \as{SG} is required only for single predictors ($u$ with one
nonzero entry), then
\begin{equation}\label{eq:dist-iii}
  \P^*\Big[\,\Pi\Big(d_H(f,f^*)>\sqrt{M\log n\cdot\log(\praw\vee n)/n}\ \Big|\ D_n\Big)>n^{-cM}\Big]\ \le\ n^{-cM},
\end{equation}
and~\eqref{eq:dist-ii} holds with $s+C_1M\log n$ in place of $s+C_1M$.
(iii) \emph{Examples.} Jointly Gaussian predictors satisfy \as{SG} with $L=1$; independent predictors satisfy it when each
satisfies it with a common constant $L$. \as{NE} holds for all large $n$ whenever every $q_n\times q_n$ principal block of
the correlation matrix has smallest eigenvalue at least a fixed $\psi>0$, for instance correlation $\rho^{\abs{j-k}}$ with
$\abs{\rho}<1$, or equal correlation $\rho\in[0,1)$.
(iv) \emph{Why each assumption.} \as{SG} is used only to make sample variances and covariances of small groups of
predictors track their population values with probability $1-n^{-cM}$. \as{NE} keeps in force the prior's penalty for
each extra predictor (a factor of order $n^{-1/2}$); this gives the rate $\rn$, the level $n^{-cM}$ and the bounded model
size. Whether \as{NE} is needed for the rate $\rn$ is open. \as{BS} bounds the signal in units of the noise.
\end{remark}

\status{Remark~\ref{rem:dist}(i)--(ii): \textbf{proved and independently checked} (same result as
Theorem~\ref{thm:dist}).}

The logarithm in the rate disappears when the true predictors are detectable.

\begin{assumption}{DS}{Detectable signal}
(a) $\lmin(\Corr(x_T))\ge\psi$ for a constant $\psi>0$ and every set $T$ of at most $(2A+3)s+2$ predictors; and (b)
$\min\{\abs{\beta^*_j}\sd(x_j)/\sigs:\beta^*_j\neq0\}\ge b_n$, where $b_n>0$ and $b_n^2n/\log n\to\infty$.
\rulesout{collinearity among small sets of predictors, and true signals below the detection scale $\rn$.}
\end{assumption}

\begin{theorem}[Parametric rate when the true predictors are detectable]\label{thm:dist-param}
Assume \as{IID}, \as{SD}, \as{SG}, \as{NE} and \as{BS} with the window $q_n$ of $M=1$, with exactly $s\ge1$ nonzero
coordinates of $\bstar$, and assume \as{DS}. Then for every sequence $M_n\to\infty$,
\[
  \Pi\big(d_H(f,f^*)>M_n/\sqrt n\ \big|\ D_n\big)\to0\quad\text{in $\P^*$-probability.}
\]
\end{theorem}

\status{\textbf{Proved and checked, conditional on} Theorem~\ref{thm:post-fixed} (in-probability form), which is itself
proved and independently checked. Verbatim: Appendix~\ref{app:status}, row 2.1.}

\paragraph{What a run returns.} The method explores the models with a sequential Monte Carlo sampler; it does not
compute the full posterior. (Small problems that it solves by exact enumeration are not covered here.) A run is
\emph{completed} when the method reports the status ``completed''; runs that stop early are not covered. The output of a
completed run contains a list $\mathcal T$ of at most $40p+1$ distinct models: the best-scoring models among those
offered to the final run of the sampler (the models scored in that run, together with models scored earlier and carried
into it, namely the precomputed tables for the blocks and the models kept from the calibration stage). Each model in
the list comes with the logarithm of its unnormalised posterior weight. The output contains no draws of
$(\beta,\sigma^2)$. Let $\pi(\mathcal T)$ be the posterior probability of the list, and let $\hat\Pi$ (not computed by
the method) be the posterior restricted to the list: the models in $\mathcal T$, weighted in proportion to the
exponentials of their returned log-weights, each with its exact conditional posterior of $(\beta,\sigma^2)$.

\inwords{If a completed run's list carries a non-negligible share of the posterior, then the posterior restricted to that
list is as good as the full posterior: it, too, puts almost no mass outside the ball of radius $\sqrt{M\log n/n}$.}

\begin{corollary}[The output of the sampler]\label{cor:dist-output}
Assume the conditions of Remark~\ref{rem:dist}(i) and that the returned log-weights are computed in exact arithmetic.
With $\P^*$-probability at least $1-n^{-cM}$ over the data, every completed run, whatever its random numbers, with
$\pi(\mathcal T)\ge n^{-cM/2}$ satisfies
\[
  \hat\Pi\Big(d_H(f,f^*)>\sqrt{M\log n/n}\Big)\ \le\ n^{-cM/2}.
\]
If a completed run scored the true model $\{j:\beta^*_j\neq0\}$ and the posterior probability of that model exceeds
$\tfrac12$, then the list contains it, it is the model the method reports as most probable, and $\pi(\mathcal T)>\tfrac12$.
\end{corollary}

\status{\textbf{Proved and checked, under exact arithmetic}. The description of the output was checked against the code
by independent reviewers. Verbatim: Appendix~\ref{app:status}, row 2.1.}

\begin{remark}[What Corollary~\ref{cor:dist-output} does not say]\label{rem:dist-output}
Whether the list carries enough posterior probability is a property of the run: the corollary does not assert it, and
the output does not certify it. Theorem~\ref{thm:dist} and Corollary~\ref{cor:dist-output} concern the exact posterior
and need no further input; Theorem~\ref{thm:dist-param} also uses the posterior concentration theorem
(Theorem~\ref{thm:post-fixed}). Statements that a completed run has in fact found and kept the true model rest, in
addition, on the sampler results of Section~\ref{sec:delivery} and on exact arithmetic.
\end{remark}

\begin{remark}[Fixed designs and the in-sample distance]\label{rem:dist-fixed}
Let the design be fixed, and let $\bar h_n$ be the in-sample average Hellinger distance between the fitted and the true
laws of one observation, on the centred scale. Assume \as{GL}, $\sz\ge1$ fixed and $2\le p\le n^A$, the energy bound
$\norm{\mu}^2\le nM$, and the floor $\lmin(\Rhat_{\Sz})\ge\varphi_0^2$ on the true model; conclusions are in probability.
(a) If also $1-\pis{\Sz}\to0$ in probability, then $\Pi(\bar h_n>M_n/\sqrt n\mid y)\to0$ for every $M_n\to\infty$
(distance, rate $M_n/\sqrt n$). (b) Without beta-min, under the superset Occam-sum condition of
Remark~\ref{rem:post-fixed}(vi) instead, $\Pi(\bar h_n^2>C'_H\log n/n\mid y)\to0$ (squared distance $\log n/n$; distance
radius $\sqrt{C'_H\log n/n}$). This covers fixed designs and the in-sample distance, which Theorem~\ref{thm:dist} does
not state.
\end{remark}

\status{\textbf{Proved and checked} (its inputs, Theorem~\ref{thm:post-fixed} and the random-design transfer lemma, are
checked). Verbatim: Appendix~\ref{app:status}, row 2.2.}

\begin{remark}[Comparison with \citet{SongLiang2015}, Theorems 1--2]\label{rem:SL-dist}
Page and equation numbers are those of the published paper.
\emph{The analogue is display~\eqref{eq:dist-i}.} Their distance is
$d(f_1,f_2)=\{\int\!\!\int(f_1^{1/2}-f_2^{1/2})^2\,dx\,dy\}^{1/2}$ between joint densities of $(x,y)$ that share the law
of $x$ (\S2.1.2, p.~952), so $d^2=2d_H^2$: the same distance. Their display~(15) (Theorem~1, p.~952; Theorem~2 states the
same form, p.~953) reads $P^*[\pi\{d(f,f^*)>\varepsilon_n\mid D_n\}>\exp(-c_1n\varepsilon_n^2)]<\exp(-c_1n\varepsilon_n^2)$;
display~\eqref{eq:dist-i} has this form with $n\varepsilon_n^2=M\log n$ (and $d_H$ in place of $d$).
\emph{Stronger.} (1) The prior is the one the method uses, with no tuning. Their theorems assume independent
Bernoulli$(r_n/P_n)$ inclusion with a hard cap $\abs{\xi}\le\bar r_n$ on the model size and conditions (9)--(14)
(pp.~951--953), which the method's beta-binomial prior with cap $\min(p,\lfloor n/2\rfloor)$ does not meet. (2) The radius
is a constant times $\rn$. At polynomial $p$ their Theorem~2 (p.~953) needs $\bar r_n\log n=o(n\varepsilon_n^2)$ with
$\bar r_n\ge r_n\to\infty$ (and their (14)), so $n\varepsilon_n^2/\log n\to\infty$; their Corollary~1 (p.~954) gives
$\varepsilon_n=O\{n^{-(1-\alpha)/2}\log^{k/2}n\}$. In display~\eqref{eq:dist-i} the radius constant does not depend on $s$,
and the smallest level $M_0$ grows at most linearly in $s$ (in the general form of the theorem, where the in-sample
signal energy and coefficient cost are bounded by fixed constants; under \as{BS} those bounds themselves grow with $s$).
(3) Theorem~\ref{thm:dist-param} gives the parametric rate $M_n/\sqrt n$ and Corollary~\ref{cor:dist-output} covers the
sampler's output; their results concern the exact posterior only and give no parametric rate. (4) No lower bound on the
growth of $p$ is needed; they assume $P_n\succ n^\theta$ (p.~950).
\emph{Weaker.} (1) $\praw\le n^A$; their Corollary~1 allows $P_n<\exp(Cn^\alpha)$. (2) Exact sparsity; their Theorem~1
covers approximate sparsity (their (3): $\lim_n\sum_j\abs{\beta^*_j}<\infty$, with predictors bounded by 1), under which the
number of nonzero coefficients may grow. (3) $s$ is fixed, with constants depending on $s$; their Theorem~2 also has a
true model of bounded size (their (4), p.~950), so this is weaker only against their Theorem~1. (4) \as{NE} is a
condition on the design that their Theorems~1--2 do not have (their eigenvalue conditions, (11) and (14), are on the
prior covariance).
\emph{Not the analogue: display~\eqref{eq:dist-iii}}, the fallback without \as{NE}. Its radius is larger by
$\sqrt{\log(\praw\vee n)}$, so its exponent is $c\,n\varepsilon_n^2/\log(\praw\vee n)$, not $c\,n\varepsilon_n^2$; and, as
proved, in the general form of the theorem (in-sample signal energy and coefficient cost bounded by fixed constants),
its squared radius at the smallest admissible level is at most of order $(1+s)^3\log n\log(\praw\vee n)/n$ for fixed $s$
(an upper bound; the proof is not known to be tight in $s$); under \as{BS} those bounds themselves grow with $s$.
\emph{Different, not ranked:} tail conditions on the predictors, \as{SG} and \as{BS} here; theirs are $\abs{x_i}\le1$
(Theorem~1), Condition~1 (an exponential-moment bound on averages of at most $\bar r_n$ predictors, Theorem~2, p.~953)
and Condition~2 (second moments, for their limit form (16), p.~954).
\emph{Overall: incomparable}: stronger on the prior, the rate and the coded output; weaker on the range of $p$, on
approximate sparsity and, for display~\eqref{eq:dist-i}, by \as{NE}.
\end{remark}

% =====================================================================================================
\section{The method, as coded, delivers the posterior}\label{sec:delivery}

\emph{The previous sections are about the posterior. This section is about the software: what the method actually
reports. It shows, in order, that the method's start law already concentrates on the true model; that the reported
model equals the true model on one data event of probability tending to one, for every seed; that a certified small
model rarely misses a true predictor, under much weaker conditions; and that every printed estimate is consistent.}

The results in this section, except Theorem~\ref{thm:screen} and Theorem~\ref{thm:screen-margin}, place their
conditions on the population law of Gaussian rows, or on sample quantities of a fixed design, that say how the signal
looks to a method that inspects blocks of at most $16$ predictors. We first define those quantities.

\paragraph{Signal gain and residual variance (fixed design).} For $U\subseteq\set{p}$ and $l\notin U$, the \emph{signal
gain} of $l$ given $U$ is $\delta^2_{U,l}:=\norm{(P_{U\cup\{l\}}-P_U)\mu}^2$: the increase in explained mean signal when
$x_l$ joins $X_U$. The \emph{residual variance} left by $U$ is $\tau^2_U:=\sigs^2+\norm{(I-P_U)\mu}^2/(n-1)$: noise plus
unexplained signal energy per observation. The sample regression coefficient of $x_l$ on $X_U$ is
$\hat\beta_{U,l}:=(X_U^\top X_U)^{-1}X_U^\top x_l$, padded with zeros to a vector indexed by a larger set when needed.

\paragraph{Partial covariances (random design).} With $W\sim\Normal_p(0,\Sig)$, $\Cov(W_j,W_l\mid W_U)$ is the partial
covariance of $W_j$ and $W_l$ given $(W_k)_{k\in U}$; $\bar\beta_{U,l}:=\Sigma_{n,UU}^{-1}\Sigma_{n,Ul}$ is the population
regression coefficient of $W_l$ on $W_U$ (padded with zeros), and $\Var(W_l\mid W_U)$ the conditional variance.

\paragraph{Leakage set and lattice.} A \emph{leakage set} is a deterministic set $Q$ with $\Sz\subseteq Q\subseteq\set{p}$
that contains every predictor able to imitate part of the signal. Its \emph{lattice} is
$\Lat(Q):=\{T:\Sz\subseteq T\subseteq Q\}$. ``No leakage'' means $Q=\Sz$.

\paragraph{Visibility.} A true predictor is \emph{visible} if its signal cannot be explained away by any $15$ other
predictors ($15$ is the largest block size minus one).

\paragraph{The signal convention.} Theorems~\ref{thm:start}, \ref{thm:reported-random} and~\ref{thm:outputs} take the
signal fixed on the standardised scale:

\begin{assumption}{FT}{Fixed standardised signals}
$\yr=\alpha\one+X\tstar+\varepsilon$ with a deterministic $\tstar$; $\sz\ge1$ is fixed; and
$\theta_{\mathrm{lo}}\le\abs{\theta^*_j}\le\theta_{\mathrm{hi}}$ for $j\in\Sz$, with fixed
$0<\theta_{\mathrm{lo}}\le\theta_{\mathrm{hi}}<\infty$.
\rulesout{true standardised coefficients that shrink to zero or blow up with $n$.}
\end{assumption}

\subsection{The start law already concentrates}\label{sec:start}

\begin{assumption}{SF}{Small-set floor}
$\lmin(\Sigma_{n,S})\ge\varphi^2$ for every $S$ with $\abs{S}\le16+\sz$, for a fixed $\varphi>0$.
\rulesout{population collinearity within sets of the size of a block; it keeps block regressions well posed.}
\end{assumption}

\begin{assumption}{V}{Visibility}
For every $j\in\Sz$ and every $U\not\ni j$ with $\abs{U}\le15$, $\abs{\Cov(W_j,W^\top\tstar\mid W_U)}\ge a$, for a fixed
$a>0$.
\rulesout{masking: a true predictor whose signal some $15$ other predictors can explain away.}
\end{assumption}

\begin{assumption}{NL}{No leakage}
For every $l\notin\Sz$ there is a set $H_l\supseteq\{l\}\cup\Sz$ with $\abs{H_l}\le e^{\kappa\sqrt{\log n}}$ for a fixed
$\kappa$, such that, for every $U\not\ni l$ with $\abs{U}\le15$ and $U_0:=U\cap H_l$:
$\abs{\Cov(W_l,W_j\mid W_{U_0})}\le Cn^{-1/2}$ for $j\in\Sz\setminus U_0$; and
$\abs{\Cov(W_l,W_v\mid W_{U_0})}$ and $\abs{\Cov(W_j,W_v\mid W_{U_0})}$ are at most $r'_n$ for $v\in U\setminus H_l$ and
$j\in\Sz\setminus U_0$, where $r'_n(n\log n)^{1/4}\to0$.
\rulesout{null predictors that imitate a true one within the start law's block view, and dependence that is not local.}
\end{assumption}

\begin{assumption}{UN}{Uncorrelated nulls}
Every null predictor is uncorrelated with every true predictor: $\Cov(W_l,W_j)=0$ for $l\notin\Sz$, $j\in\Sz$.
\end{assumption}

\inwords{The start law looks at the predictors one block at a time. If every true predictor is visible and no null
predictor leaks signal, the start law already puts mass $1-n^{-c}$ on the true model, for every $c<\frac12$. If, in
addition, the null predictors are uncorrelated with the true ones, the start law's miss is $\OP(\rn)$, as small as the
posterior's up to a factor $\sqrt{\log n}$.}

\begin{theorem}[The start law concentrates]\label{thm:start}
Assume \as{GR}, \as{GL}, \as{DIM} with $p>\sz$, \as{FT}, \as{SF}, \as{V} and \as{NL}. Then, for every $c<\tfrac12$,
\[
  \P\big(\pi_0(\Sz)\ge1-n^{-c}\big)\to1 .
\]
If moreover \as{UN} holds, then $1-\pi_0(\Sz)=\OP(\rn)$. The exponent $\tfrac12$ cannot be improved:
$1-\pi_0(\Sz)\ge\tfrac13(n+1)^{-1/2}$ for every valid data set with $n\ge4$ and $\sz<K$ (hence for all large $n$). So the
rate under \as{UN} is optimal up to the factor $\sqrt{\log n}$.
\end{theorem}

\status{\textbf{Proved and independently checked}. Verbatim: Appendix~\ref{app:status}, row 3.3.}

\begin{remark}[Leakage, and why each assumption]\label{rem:start}
(i) \emph{With leakage.} If \as{NL} holds with $\Sz$ replaced by a deterministic $Q\supseteq\Sz$ with $\abs{Q}\le K$
(predictors outside $Q$ do not leak), then
$\P\big(\pi_0(\Lat(Q))\ge1-(\abs{Q}+1)n^{-1/2}e^{C\sqrt{\log n}}\big)\to1$; for $\abs{Q}\le n^{1/2-\epsilon}$ this is at
least $1-n^{-c}$ for every $c<\epsilon$.
(ii) \as{V} says that each true predictor, alone, improves the fit of every small model it is missing from; then the
start law, which sees only within-block interactions, still ranks $\Sz$ first. The level is sharp in the fixed-design
analogue (threshold $\log n+2\log p$ for every $p$).
(iii) In \as{NL} the partial-covariance level $n^{-1/2}$ is a threshold of the proof (numerics show no cliff there);
localisation is needed because the blocks depend on the data. The size bound on $H_l$ only counts reference sets in a
union bound; it is not known to be needed.
(iv) \as{UN} is an additional hypothesis for the sharper rate only; it does not replace \as{NL}. It suffices that these
correlations be at most $C/\sqrt{n\log n}$ given every proper subset of the truth. Not known to be necessary.
(v) This theorem is about an internal object of the method. It explains why the sampler starts in the right place.
\end{remark}

\subsection{The reported model is the true model}\label{sec:reported}

\emph{This is the main result: on one data event whose probability tends to one, every run that produces output reports
exactly the true model, whatever its random seed. We state it for Gaussian rows, with conditions on the population law,
and then for fixed designs, with conditions on the sample.}

\paragraph{Random Gaussian design.} Fix $b\in(0,\tfrac12)$ and $\rho>0$ with $15\rho+b<\tfrac12$.

\begin{assumption}{PL}{Polynomial leakage set}
There is a deterministic $Q\supseteq\Sz$ with $\abs{Q}\le n^b$ and, for large $n$, $\abs{Q}\le K$ and ($Q=\set{p}$ or
$K\ge2\abs{Q}+32\sz+1$), such that:
(i) \emph{(no leak outside $Q$)} there is a deterministic $\varpi_n$ with $\varpi_n\sqrt{n/\log n}\to0$ such that
$\abs{\Cov(W_l,W_j\mid W_U)}\le\varpi_n$ for all $l\notin Q$, $j\in\Sz\setminus U$ and $\abs{U}\le15$;
(ii) \emph{(locality)} there is a deterministic $r'_n\to0$, and every $l\notin Q$ has a deterministic
$H_l\supseteq\{l\}\cup\Sz$ with $\abs{H_l}\le n^{\rho}$, such that for all $l\notin Q$, all $U\not\ni l$ with
$\abs{U}\le15$ and $U_0:=U\cap H_l$: $\abs{\Cov(W_l,W_v\mid W_{U_0})}\le r'_n$ and $\abs{\Cov(W_j,W_v\mid W_{U_0})}\le r'_n$
for all $v\in U\setminus H_l$ and $j\in\Sz\setminus U_0$.
\rulesout{many predictors that can imitate the signal: those that can must form a set of size at most $n^b$, $b<\tfrac12$.}
\end{assumption}

\begin{assumption}{PQ}{Floor on the leakage set}
$\lmin(\Sigma_{n,Q\cup\{l\}})\ge\varphi^2$ for every $l\in\set{p}$, with $\varphi$ as in \as{PE}.
\rulesout{population collinearity between the leakage set and any other predictor.}
\end{assumption}

\begin{assumption}{LR}{Local regressions within the leakage set}
There are deterministic sequences $\nu_n=o(\log n)$ and $L_n$, and for every $l\in\set{p}$ a deterministic
$N_l\subseteq Q\setminus\{l\}$ with $\abs{N_l}\le\nu_n$, such that
$\norm{\bar\beta_{T,l}-\bar\beta_{T\cap N_l,l}}_1\le L_n$ for every $T$ with $\Sz\subseteq T\subseteq Q$ and $l\notin T$,
where $L_n\to0$ and $L_n^2\log(\abs{Q}+1)=o(\log n)$.
\rulesout{regressions, inside the leakage set, that depend on many members of it at once.}
\end{assumption}

\inwords{If every true predictor is visible, the predictors that can imitate the signal form a set $Q$ of size at most
$n^b$ with $b<\tfrac12$, and regressions within $Q$ are local, then, with probability tending to one, every run that
produces output reports the true model. The start mode contains every true predictor and possibly some members of $Q$;
the local search then removes the spurious ones, one per move.}

\begin{theorem}[The reported model equals the truth: random Gaussian design]\label{thm:reported-random}
Let $\Sig$ be positive definite and assume \as{GR}, \as{GL}, \as{DIM}, \as{MC}, \as{FT}, \as{V}, \as{PL}, \as{PE},
\as{PQ} and \as{LR}. Then, under the joint law of the design and the noise, there are data events $\mathcal G_n$ with
$\P(\mathcal G_n)\to1$ on which \emph{every} admissible run that produces output reports $\Shat=\Sz$, whatever its seed
and the outcome of its internal calibration. In particular $\P(\text{output and }\Shat\neq\Sz)\to0$.

If moreover \as{DC} holds, then $\P(\text{the run produces output})=1$ for every admissible $n\ge3$ with $(n,\praw)$ in the
bounded domain $\Dom$ on which the code's two fixed integer limits are not reached. If in addition (DC$^+$) holds,
\[
  \P(\text{the run produces output})=1\ \ \text{for every admissible }n\ge3,\qquad \P(\Shat=\Sz)\to1 .
\]
\end{theorem}

\status{\textbf{Proved and checked}; its one condition (the transfer from population to sample conditions) is itself
proved and independently checked. Delivery clause: proved and independently checked. \bridgepending{} A revision of this statement from a completed
strengthening is pending. Verbatim: Appendix~\ref{app:status}, row 3.1.}

The two integer limits are a base search cap of $10^6$ and a 32-bit integer bound $2^{31}-1$ on the search cap. The first
can bind only when $n+p+1$ exceeds $2^{10^6}$, the second only when the starting model has more than about
$2.1\times10^9$ predictors, which needs $n\gtrsim4.3\times10^9$ since $\abs{\tSz}\le\lfloor n/2\rfloor$. So
\[
  \Dom=\big\{(n,\praw):\lceil\log_2(n+\praw+1)\rceil\le10^6,\ \lceil\log_2(n+\praw+1)\rceil+\lfloor n/2\rfloor\le2^{31}-1\big\}.
\]

\begin{remark}[On Theorem~\ref{thm:reported-random}]\label{rem:reported-random}
(i) On $\mathcal G_n$ the local search from the start mode only deletes the spurious members of the start mode, one per
move, in at most $n^b$ moves, and certifies $\Sz$; its move cap never binds.
(ii) The case $\abs{Q}\le(\log n)^\kappa$ is a special case.
(iii) Examples, never hypotheses: correlation that is local in a metric, polynomial correlation decay of exponent
$\alpha>2$ (not yet checked), and one-factor blocks.
(iv) \emph{Why $b<\tfrac12$.} At model size about $\sqrt n$ the prior-plus-Occam penalty of one more predictor
($\approx\log p+\tfrac12\log n$) stops dominating the best of $p$ noise gains ($\approx\log p$); beyond it the key events of
the argument are expected to fail (heuristic; whether the conclusion fails is open). \as{LR} is a sufficient form of a
width condition that the argument needs; it is not known to be necessary for the conclusion.
(v) All assumptions are on the population law. \as{PL} and \as{LR} are long and specific; they are the price of a
statement about a \emph{search} algorithm: the signal is not masked, what can imitate it is confined to a set of size
$n^b$, and within that set regressions are local.
\end{remark}

\begin{remark}[Comparison with the literature]\label{rem:SL-reported}
This is consistency of what the software prints, with $p\gg n$ and no design family. We know of no comparable result
for the selection output of an SMC or MCMC implementation of Bayesian variable selection: \citet{SongLiang2015} run the
SAMC sampler of \citet{LiangLiuCarroll2007}, citing its general properties, but prove nothing about its selection
output.
\end{remark}

\paragraph{Fixed design.} The same conclusion holds for a fixed design, with non-constant raw columns, under sample-level
conditions.

\begin{assumption}{FF}{Floors for blocks and the leakage set}
$\lmin(\Rhat_T)\ge\varphi_0^2$ for every $T$ with $\abs{T}\le16$ and for every $T=Q\cup\{l\}$, $l\in\set{p}$, with $Q$ as
in \as{BL}.
\rulesout{collinearity within a block, or between the leakage set and any other predictor.}
\end{assumption}

\begin{assumption}{VF}{Visibility, fixed design}
There is a deterministic $M_n\to\infty$ such that $\delta^2_{U,j}\ge M_n\log n$ for every $j\in\Sz$ and every $U\not\ni j$
with $\abs{U}\le15$.
\rulesout{masking in the sample: a true predictor whose signal some $15$ others explain away.}
\end{assumption}

\begin{assumption}{BL}{Bounded leakage}
There are a deterministic $Q\supseteq\Sz$ (with $\abs{Q}\le K$ and ($Q=\set{p}$ or $K\ge2\abs{Q}+32\sz+1$) for large $n$)
and a fixed $\eta\in[0,\eta_A)$, $\eta_A:=(\sqrt{2A+1}-\sqrt{2A})^2$, such that $\delta^2_{U,l}\le\eta\,\tau^2_U\log n$ for
every $l\notin Q$ and every $U\not\ni l$ with $\abs{U}\le15$.
\rulesout{predictors outside $Q$ that, added to a small model, explain more than a fraction $\eta$ of $\log n$
residual-variance units.}
\end{assumption}

\begin{assumption}{LQ}{Size of the leakage set}
One of: (a) $\abs{Q\setminus\Sz}\le(1-\delta)\log n$ for a fixed $\delta\in(0,1)$; or (b) $\abs{Q}\le n^b$ with
$b<\tfrac12-\tfrac\eta2-\sqrt{2A\eta}$, $\theta^{*2}_{\min}n/(\abs{Q}+\log n)\to\infty$, and local regressions: for every $l$
a deterministic set $N_l\subseteq Q\setminus\{l\}$ with $\abs{N_l}\le h\log_2n$, $h+b<\tfrac12$, and a deterministic
$L_n\to0$ with $L_n^2\log(\abs{Q}+1)=o(\log n)$, such that $\norm{\hat\beta_{T,l}-\hat\beta_{T\cap N_l,l}}_1\le L_n$ for every
$T$ with $\Sz\subseteq T\subseteq Q$, $\abs{T}<K$, $l\notin T$.
\rulesout{a leakage set so large, or so entangled, that adding a predictor could improve the score somewhere between
$\Sz$ and $Q$.}
\end{assumption}

\begin{theorem}[The reported model equals the truth: fixed design]\label{thm:reported-fixed}
Let the design be fixed, with non-constant raw columns. Assume \as{GL}, \as{DIM}, \as{MC}; \as{SP}, \as{BC}, \as{SE},
\as{NC} and \as{BM}(b); and \as{FF}, \as{VF}, \as{BL} and \as{LQ}. Then the conclusions of
Theorem~\ref{thm:reported-random} hold: on data events $\mathcal G_n$ with $\P(\mathcal G_n)\to1$, every admissible run
that produces output reports $\Shat=\Sz$, for every seed and calibration outcome, and the delivery clauses hold under
\as{DC} on $\Dom$ and under (DC$^+$) for every $n$.
\end{theorem}

\status{\textbf{Proved and independently checked}. \bridgepending{} A revision from a completed strengthening is pending. Verbatim: Appendix~\ref{app:status},
row 3.2.}

\begin{remark}[Intuition and limits]\label{rem:reported-fixed}
\as{VF}: each true predictor, added alone to any small model that misses it, explains much more than $\log n$ of the
signal, more than the prior charges for one predictor; so the start law never leaves a truth out. \as{BL}: outside $Q$,
no predictor looks like signal, so none enters the start mode. \as{LQ}: predictors in $Q$ may enter the start mode; the
local search removes them one at a time, which needs that the score never rewards \emph{adding} a predictor between $\Sz$
and $Q$; this holds when $Q$ has fewer than $\log n$ extra predictors (a), or when regressions on such models are carried
by a few members of $Q$ (b). These are sample-level conditions on $(X,\tstar)$; they are interpretable but not reducible
to eigenvalue floors and beta-min, since visibility fails on well-conditioned designs where truths cancel.
Theorem~\ref{thm:reported-random} shows that population analogues imply them with probability tending to one. Not
available: the theorem without any condition on the lattice (control of the actual search path), and any $b\ge\tfrac12$.
\end{remark}

\subsection{Certified screening}\label{sec:screen}

\emph{The previous theorems need visibility and bounded leakage. The next result needs neither: it says only that when
the method certifies a small model, that model rarely misses a true predictor. Its conditions can hold on designs, such
as designs with strong factors, where those of Section~\ref{sec:reported} fail.}

In this subsection $q_n$ denotes a model-size budget chosen by the analyst (it is not the window of
Section~\ref{sec:distribution}), and $M$ a bound on the coefficients.

\begin{assumption}{EK}{Sparse eigenvalue floor $\kappa$}
$\lmin(\Rhat_S)\ge\kappa$ for every $S$ with $\abs{S}\le q_n+\sz$, with a constant $\kappa>0$ fixed in $n$.
\rulesout{collinearity among small sets of predictors.}
\end{assumption}

\begin{assumption}{FC}{Signals and coefficients bounded by fixed constants}
$\tmin\ge b_0>0$ and $\norm{\tstar}_1\le M$, where $\kappa$, $b_0$, $M$ and $\sigs^2$ do not depend on $n$.
\rulesout{signals that vary with $n$ (see Theorem~\ref{thm:screen-margin}).}
\end{assumption}

\begin{assumption}{SB}{Size budget}
$(q_n+1)\log p+\log n\le c_{\mathrm{SB}}\,n$, where $c_{\mathrm{SB}}:=\kappa^2b_0^2/\{200(M^2+\sigs^2)\}$ and $q_n\in\{0,\dots,K-1\}$.
\rulesout{a model-size budget $q_n$ too large for a union bound over models of that size.}
\end{assumption}

\inwords{If the design has a sparse eigenvalue floor, the signals and coefficients are bounded by fixed constants, and
$q_n$ is below the model-size cap, then the probability that a run produces output, certifies a reported model of size
at most $q_n$, and that model misses a true predictor tends to zero, with no visibility, leakage, start or block
condition.}

\begin{theorem}[Certified screening]\label{thm:screen}
Let the design be fixed and assume \as{GL}, \as{MC}, \as{EK}, \as{FC} and \as{SB}. Then
\[
  \P\big(\text{output},\ \Shat\text{ certified},\ \abs{\Shat}\le q_n,\ \Sz\not\subseteq\Shat\big)\to0 .
\]
In particular this holds for $p\le n^A$ and every $q_n=o(n/\log n)$.
\end{theorem}

\status{\textbf{Proved and independently checked}. Verbatim: Appendix~\ref{app:status}, row 3.7.}

No condition on the start, leakage, blocks or visibility is used, and $p\le n^A$ is not needed. The argument needs only a
weaker, directional floor, which leaves correlations among null predictors free; the full floor \as{EK} does constrain
them.

For signals that may vary with $n$, define the \emph{noiseless gain} $\gamma(t\mid S)$ of $t\notin S$ given $S$: the
decrease of the method's ridge-regularised residual quadratic form when $t$ joins $S$ and the response is replaced by its
mean $\mu$. Let
$\Gamma_n(q):=\min\{\max_{t\in\Sz\setminus S}\gamma(t\mid S):\abs{S}\le q,\ S\not\supseteq\Sz\}$, so that every model of
size at most $q$ that misses a truth has some missing truth with gain at least $\Gamma_n(q)$, and let
$\tau_n^2:=\sigs^2+\norm{\mu}^2/(n-1)$.

\begin{assumption}{GM}{Gain margin}
$q_n\in\{0,\dots,K-1\}$ and
\[
  \liminf_{n\to\infty}\frac{\sqrt{\Gamma_n(q_n)}}{\tau_n\sqrt{2\log p+\log n}+\sigs\sqrt{2q_n\log p}}>1 .
\]
\rulesout{small models that miss a truth while every missing truth has a gain below the noise level of a union bound.}
\end{assumption}

\begin{theorem}[Certified screening with signals that may vary with $n$]\label{thm:screen-margin}
Let the design be fixed and assume \as{GL}, \as{MC} and \as{GM}. Then
$\P\big(\text{output},\ \Shat\text{ certified},\ \abs{\Shat}\le q_n,\ \Sz\not\subseteq\Shat\big)\to0$. No condition
$p\le n^A$, no bound on $\sz$, and no condition on the start, leakage, blocks or the collinearity of the null predictors is
used.
\end{theorem}

\status{\textbf{Proved and independently checked} (a non-asymptotic form, with probability at most $3e^{-x}$ under a
finite-$n$ margin, is also checked). Verbatim: Appendix~\ref{app:status}, row 3.7.}

\begin{remark}[What certified screening does not say]\label{rem:screen}
(i) The bound is on the probability of an intersection: a run produces output, certifies a small model, \emph{and} that
model misses a true predictor. It is not a guarantee conditional on certification, and it does not say that the method
returns a superset of the true model with probability tending to one.
(ii) A run is not claimed to certify a small model with probability tending to one; how often it does is reported from
numerics below.
(iii) It is one-sided: certified small models are not shown to equal the true model. Exact recovery from the certificate
needs a near-orthogonality floor (proved and checked), which is not reader-friendly.
\end{remark}

\paragraph{How often runs certify (numerics, not a theorem).} In existing numerics the method certified the reported
model in every completed run: 2,394 of 2,394 in the 2,400-fit study ($n=200$); 376 of 376 in the strong-factor experiment
($n\le150$, $p$ up to $10n$); and 71 of 71 in the screening study ($n=100$). Where model sizes were recorded, the certified
models were small: all 376 strong-factor fits had at most $2\sz$ predictors, and all 71 screening fits at most 5 (the
2,400-fit study records no per-fit size). On strong-factor designs (equicorrelation $\rho\in\{0.5,0.7\}$, $p\ge2.5n$,
$\sz\in\{4,8\}$) the certified model contained every true predictor in 61 of 122 fits at $n=50$, 18 of 20 at $n=75$ and 30
of 30 at $n=100$. At these constants \as{SB} requires $n\gtrsim1{,}450$.

\subsection{The printed outputs are consistent}\label{sec:outputs}

\emph{The method prints an evidence estimate, pooled inclusion probabilities and the corresponding median probability
model. This subsection shows that all of them are consistent, on the event that the method's calibration succeeds.}

This result concerns the \emph{sampler branch} of the method, used when the model space is too large to enumerate; the
exact-enumeration branch is not covered. Let $\Acal$ be the event that the method takes the sampler branch and its
calibration succeeds (fallbacks included); every completed sampler run lies in $\Acal$. On $\Acal$, let $\Zhat$,
$\PIPhat$ and $\MPMhat$ be the evidence estimate, the pooled inclusion probabilities and their median probability model,
computed from the four islands' uninterrupted records with equal weights (Section~\ref{sec:method}). On completed runs
they are exactly the printed fields. Let $\rho_n:=n^{-1/2}e^{C\sqrt{\log n}}$ for a suitable constant $C$; so
$\rho_n\le n^{-c}$ eventually, for every $c<\tfrac12$.

\inwords{Under the conditions that make the start law concentrate, every printed estimate is within $\rho_n$ of its
target, on the event that the calibration succeeds: the evidence estimate is close to the evidence, the inclusion
probabilities are close to $0$ or $1$, and the printed median probability model and reported model are the true model.}

\begin{theorem}[The printed outputs are consistent]\label{thm:outputs}
Assume \as{GR}, \as{GL}, \as{DIM} with $p>\sz$, \as{MC}, \as{SC}, \as{FT}, \as{PE}, \as{V} and \as{NL}. Then, under the
joint law of the design, the noise and the method's random numbers,
\[
  \ind{\Acal}\Big(\big|\Zhat/Z-1\big|+\big|\log\Zhat-\log Z\big|+\max_{j}\big|\PIPhat_j-\ind{j\in\Sz}\big|\Big)=\OP(\rho_n),
\]
\[
  \P\big(\Acal\cap\{\MPMhat\neq\Sz\}\big)\to0,\qquad \P\big(\text{the run completes and }\Shat\neq\Sz\big)\to0,
\]
and the same rate holds for the printed fields multiplied by the indicator of a completed run.
\end{theorem}

\status{\textbf{Proved and independently checked}. \bridgepending{} Verbatim: Appendix~\ref{app:status}, row 3.6.}

\begin{remark}[On Theorem~\ref{thm:outputs}]\label{rem:outputs}
(i) No lower bound on $\P(\Acal)$ is asserted: the statements are about intersections with $\Acal$ or with completion,
not conditional on them. Under \as{DC}, completion holds on $\Dom$, and under (DC$^+$) for every admissible $n$; that link
is not used here.
(ii) The rate $\rho_n$ is the start law's rate, slower than the posterior rate $\rn$ by a factor $e^{C\sqrt{\log n}}$.
Asymptotically, the start law concentrates on $\Sz$, so the sampler has vanishing mass to transport and $\Zhat/Z\to1$,
whatever the length of the tempering ladder. At moderate $n$ the posterior is a mixture over $\Sz$ and its many
one-predictor supersets (Remark~\ref{rem:moderate}). An alternative argument, with no condition on $\Zhat$, gives the
same conclusions for the inclusion probabilities and the median and highest posterior models under a global-overlap
condition.
\end{remark}

\begin{remark}[The pooled evidence estimate is unbiased]\label{rem:unbiased}
Given the calibration, each island's evidence estimate is unbiased for $Z$, for every design and data set; no condition
on the design is needed. By linearity, so is the printed equal-weight pooled estimate: on $\Acal$,
$\E[\Zhat\mid\text{calibration}]=Z$. This is the classical unbiasedness of sequential Monte Carlo estimates of normalising
constants \citep{DelMoral2004}, established for the method as coded, calibration included: the calibration is frozen
before the final run's random numbers are drawn.
\end{remark}

\status{Per island: \textbf{proved and independently checked}. Pooled form: \textbf{proved (a two-line argument), not
yet independently checked}. Verbatim: Appendix~\ref{app:status}, row 3.5.}

% =====================================================================================================
\section{The parallel method}\label{sec:parallel}

\emph{The method runs four independent islands and pools them. What is proved about the pooling is short. A parallel
engine that reproduces the serial output exactly inherits every guarantee above, but that inheritance is not yet
written as a theorem. Running many islands at very large $p$ raises questions that are open.}

\paragraph{Pooling over the four islands (proved).} The pooled inclusion probability is the equal-weight mean of the
island estimates, so its sup-norm error is at most that of the worst island:
$\max_j\abs{\PIPhat_j-\PIP_j}\le\max_r\max_j\abs{\PIPhat_{r,j}-\PIP_j}$. The pooled evidence estimate
$\Zhat=\frac14\sum_r\Zhat_r$ is unbiased given the calibration (Remark~\ref{rem:unbiased}).

\status{Worst-island bound: \textbf{proved and independently checked}. Pooled evidence: \textbf{proved, not yet
independently checked}. Verbatim: Appendix~\ref{app:status}, row 4.1.}

\paragraph{An output-identical parallel engine (empirical; no theorem).} The current parallel engine reproduces the serial
engine's output exactly on the tested set: byte-identical output in 24 of 24 production test cases and 10 of 10
supplementary cases. If an engine's output is identical to the serial method's, every result of
Sections~\ref{sec:selection}--\ref{sec:delivery} applies to it unchanged; a note stating this inheritance precisely is in
progress.

\status{\textbf{In progress}: empirical fact and a note in progress; no theorem.}

\paragraph{Many islands and island weighting at very large $p$ (open).} Runs at larger $p$ show that the spread of the
islands' log-evidence estimates grows with $p$: $0.09$ at $p=2{,}000$, and $0.76$--$1.15$ for $p$ between $5{,}000$ and
$20{,}000$, where one island held $79\%$ of the evidence weight at $p=5{,}000$. The islands' inclusion probabilities
nevertheless agree to within $5\times10^{-4}$. This is why the method pools islands with equal weights rather than by
their evidence estimates. Results for many islands, for weighting islands, and for a fixed number of particles per
island at $p=10^6$ are open. The planned run at $p=10^6$ uses an engineered variant of the method; it will be reported as
a variant validated against the method of this paper on shared problem sizes, not as a run covered by the theorems.

\status{\textbf{Open}; no theorem.}

% =====================================================================================================
\section{Open questions}\label{sec:open}

\emph{We list the questions our results leave open, in the order of the sections above. None of them is needed for any
statement in this paper.}

\begin{enumerate}[leftmargin=1.5em]
  \item \emph{Over-fitted models with no eigenvalue condition.} Find a sufficient condition for the superset Occam-sum
    condition of Remark~\ref{rem:post-fixed}(vi) that involves no design quantity. Whether that condition holds without
    any eigenvalue condition is open. A positive answer would remove the vanishing floor \as{NC} from
    Theorem~\ref{thm:post-margin}, so that it would need no eigenvalue condition away from the true model.
  \item \emph{The floor on the truth.} Replace the floor $\lmin(\Rhat_{\Sz})\ge\varphi_0^2$ in \as{ID} by a bounded ridge
    cost; the floor is sufficient and not known to be necessary.
  \item \emph{The margin constant.} The constant $c_1(\sz,A)$ in \as{ID} is explicit but not optimised.
  \item \emph{A lower bound on every design.} Prove $\E[1-\pis{\Sz}]\gtrsim\rn$ under a floor on the sets
    $\Sz\cup\{j\}$, matching Theorem~\ref{thm:post-fixed} on every such design (Remark~\ref{rem:sharp}(vi)).
  \item \emph{The $\sqrt{\log n}$ in probability when $\log p=o(\log n)$}, and the matching upper bound in expectation on
    raw Gaussian designs (Remark~\ref{rem:sharp}(vi)).
  \item \emph{Whether \as{NE} is needed} for the rate $\rn$ in Theorem~\ref{thm:dist}.
  \item \emph{Leakage sets of size $n^b$ with $b\ge\tfrac12$}, and a reported-model theorem with no condition on the
    lattice of the leakage set, which would require control of the actual search path
    (Remark~\ref{rem:reported-fixed}).
  \item \emph{Raw-scale signals} in Theorems~\ref{thm:start}, \ref{thm:reported-random} and~\ref{thm:outputs}, which
    are stated with the signal fixed on the standardised scale \as{FT}.
  \item \emph{Certified screening with signals that vary with $n$} under a sparse eigenvalue floor, in the plain form of
    Theorem~\ref{thm:screen}.
  \item \emph{Sharper rates for the printed outputs} (Theorem~\ref{thm:outputs}).
  \item \emph{The parallel method}: the inheritance note, and results for many islands and island weighting
    (Section~\ref{sec:parallel}).
\end{enumerate}

\bibliography{../shared/refs}

\appendix
\section{Notation}\label{app:notation}

\begin{longtable}{@{}p{0.27\linewidth}p{0.55\linewidth}l@{}}
\toprule
Symbol & Meaning & Where \\ \midrule \endhead
$n$; $p$; $\praw$ & sample size; number of retained (non-constant) predictors; raw number of predictors & \S\ref{sec:model} \\
$\yr$, $\Xr$ & raw response and predictors & \S\ref{sec:model} \\
$\alpha^*,\bstar,\sigs^2$ & true intercept, coefficients and noise variance & \S\ref{sec:model} \\
$\Sz$, $\sz$ & true model $\supp\bstar$ and its size & \S\ref{sec:model} \\
$\Sig$, $W$, $\Sigma_{n,S}$ & population correlation matrix of a row; a population row; principal sub-matrix & \S\ref{sec:model} \\
$P^\perp_{\one}$, $D$, $s_y$ & centring matrix; diagonal of column sample s.d.; response s.d. & \S\ref{sec:standardise} \\
$X$, $x_j$, $y$ & standardised design and column $j$ ($\norm{x_j}^2=n-1$); standardised response & \S\ref{sec:standardise} \\
$\tstar$, $\mu$, $\tmin$ & standardised coefficients $D\bstar$; mean signal $X\tstar$; smallest true $\abs{\theta^*_j}$ & \S\ref{sec:standardise} \\
$X_S$, $G_S$, $\Rhat_S$, $P_S$ & columns in $S$; Gram matrix; sample correlation matrix; projection & \S\ref{sec:standardise} \\
$\Om$, $K$ & model space; size cap $\min(p,\lfloor n/2\rfloor)$ & \S\ref{sec:prior} \\
$w(k)$, $\varpi$ & model-size weight $1/\{2\binom{2p-1}{k}\}$; capped beta-binomial prior & \S\ref{sec:prior} \\
$a_0,b_0$ & inverse-gamma hyperparameters ($0.01$ each) & \S\ref{sec:prior} \\
$m(S)$, $H_S$, $Q_S$, $z_S$, $\nu$ & marginal likelihood and its ingredients & \S\ref{sec:posterior} \\
$\pis{S}$, $F(S)$, $Z$ & posterior over models; log score; evidence & \S\ref{sec:posterior}, \S\ref{sec:method} \\
$\PIP_j$, $\HPM$, $\MPM$ & inclusion probability; highest posterior model; median probability model & \S\ref{sec:posterior} \\
$\miss$ & posterior miss $1-\pis{\Sz}$ & \S\ref{sec:posterior} \\
$C_b$, $J$ & blocks of the method's partition ($\abs{C_b}\le16$); number of blocks & \S\ref{sec:method} \\
$\pi_0$, $\tSz$ & start law; start mode & \S\ref{sec:method} \\
$\ell_n$ & base move cap $\lceil\log_2(n+p+1)\rceil$ & \S\ref{sec:method} \\
$\Shat$ & reported model & \S\ref{sec:method} \\
$R$; $\PIPhat_{r,j}$, $\Zhat_r$ & number of islands (4); island estimates & \S\ref{sec:method} \\
$\PIPhat_j$, $\Zhat$, $\MPMhat$ & pooled estimates and estimated median probability model & \S\ref{sec:method} \\
$\Dom$ & bounded domain of $(n,\praw)$ on which the code's integer limits are not reached & \S\ref{sec:standing}, \S\ref{sec:reported} \\
$A$, $\rate$ & dimension exponent ($p\le n^A$); reference rate $\rn$ & \S\ref{sec:standing}, \S\ref{sec:asymptotics} \\
$\OP,\oP,O,o,\lesssim,\asymp,\gg$ & asymptotic notation & \S\ref{sec:asymptotics} \\
$\varphi_0,g_0,\bar\theta,C_\beta$ & fixed constants of the fixed-design assumptions & \S\ref{sec:fixed} \\
$c_0(\sz,A)$, $c_1(\sz,A)$ & beta-min and margin constants & \S\ref{sec:fixed} \\
$\varphi$, $b_{\mathrm{lo}},b_{\mathrm{hi}}$ & population floor; bounds on true raw coefficients & \S\ref{sec:random} \\
$\mathcal C$, $\beta^{\mathcal C}$, $\sigma^2_{\mathcal C}$, $\Sz^{\mathcal C}$, $s$ & candidates; best linear predictor; residual variance; its support and size & \S\ref{sec:omitted} \\
$D_n$, $\Pi(\cdot\mid D_n)$, $\P^*$ & data set; posterior of $(S,\beta,\sigma^2)$; law of the data & \S\ref{sec:dist-setting} \\
$f$, $f^*$, $d_H$, $H^2$ & fitted and true joint densities; Hellinger distances & \S\ref{sec:dist-setting} \\
$M$, $q_n$ & level and window of the distribution theorem & \S\ref{sec:dist-setting} \\
$\mathcal T$, $\hat\Pi$ & list returned by a completed run; posterior restricted to it & \S\ref{sec:dist-results} \\
$\delta^2_{U,l}$, $\tau^2_U$, $\hat\beta_{U,l}$ & signal gain; residual variance; sample regression coefficient & \S\ref{sec:delivery} \\
$\Cov(\cdot,\cdot\mid W_U)$, $\bar\beta_{U,l}$ & partial covariance; population regression coefficient & \S\ref{sec:delivery} \\
$Q$, $\Lat(Q)$ & leakage set; its lattice & \S\ref{sec:delivery} \\
$\Acal$, $\rho_n$ & sampler branch with successful calibration; rate $n^{-1/2}e^{C\sqrt{\log n}}$ & \S\ref{sec:outputs} \\
\bottomrule
\end{longtable}

\section{Delivery and sampler conventions in full}\label{app:conventions}

\paragraph{Delivery conventions \as{DC}.} Used only by the clauses ``the run always produces output'' of
Theorems~\ref{thm:reported-random} and~\ref{thm:reported-fixed}. They are operational assumptions, not consequences of any
statistical condition.
\begin{enumerate}[label=(\roman*),leftmargin=2em]
  \item \emph{Idealised machine:} exact arithmetic at unit cost (including $\exp$ and $\log$); no integer, index,
    stream-identifier or allocation overflow; wall time at most $\tau\cdot(\text{operation count})+\tau_0$ and peak memory at
    most $\beta\cdot(\text{words held})+\beta_0$ for fixed constants $(\tau,\tau_0,\beta,\beta_0)$; every allocation within
    the declared memory succeeds; the build and its dependencies work; no external interruption or worker crash.
  \item \emph{Covered configuration:} the production core with the local search from the start mode (cap
    $\ell_n+\abs{\tSz}$) and the tie-aware start-mode fallback, implemented with
    $O\big(p\sum_b(2^{\abs{C_b}}+(K+1)^2)\big)$ operations and $O\big(\sum_b2^{\abs{C_b}}+J(K+1)\big)$ storage; the local
    search from the start mode starting at the smallest, lexicographically first maximiser of the start law; the default
    list of summaries only; certificate computations with declared polynomial operation and memory budgets, included in
    the allocation; serial execution or a fixed number of concurrent workers (memory envelope multiplied by that number);
    a fixed number of islands; the grid fallback switched on (the shipped default).
  \item \emph{Random numbers and stops:} ideal independent uniform random numbers; the method's stop-and-return rule and
    stop-only resource guard as coded (a stopped run returns no estimates).
  \item \emph{Inputs and allocation:} a valid seed (an integer in $[1,2^{31}-1]$) and a declared memory and time allocation
    at least the method's explicit polynomial work and memory envelopes in $(n,\praw)$, including the preflight bounds and
    the local-search archive and time bounds.
\end{enumerate}
Under (i)--(iv) the delivery clauses hold on the bounded domain $\Dom$ of Section~\ref{sec:reported}, on which the code's
two fixed integer limits (a run is rejected when the base cap $\ell_n$ exceeds a fixed ceiling, or when the start-mode cap
$\ell_n+\abs{\tSz}$ exceeds the integer range) are not reached. \emph{Convention (DC$^+$)} is the idealised machine of (i)
with these two limits lifted, everything else (other validators, the start-mode enumeration cap and its fallback, every
method cap) kept as coded; under it the delivery clauses hold for every $n$. Plain (i) does not give delivery for every
$n$.

\paragraph{Sampler conventions \as{SC}.} Used only by Theorem~\ref{thm:outputs} and Remark~\ref{rem:unbiased}. Exact
arithmetic; within each random stream, fresh independent ideal uniform random numbers; the final run's stream for each
island independent of the data and of the calibration's streams; the coded start sampler returns exact draws from
$\pi_0$, and the reported start normaliser equals the exact normalising constant of $\pi_0$'s unnormalised weights; the
coded Markov moves realise their stated kernels, applied in a uniformly random order independent of the state; the
calibration record is a function of the data and the calibration streams; each island's record is a measurable function
of the calibration value and its stream; statements about estimates refer to the potential (uninterrupted) run,
identified with the reported fields on completed runs.

\section{Status ledger (internal; removed for submission)}\label{app:status}

\IfFileExists{status_ledger.tex}{\begingroup\footnotesize
\begin{longtable}{@{}p{0.15\linewidth}p{0.05\linewidth}p{0.5\linewidth}p{0.22\linewidth}@{}}
\toprule Result & Card row & Registry status, verbatim (\texttt{claims.tsv}, main f8a6c2da) & Label used \\ \midrule \endhead
Theorem 1, in-probability form (HPM/MPM/PIPs) & 1.1 & \texttt{SOP:TH09@1.0.2}: "PROVED-CHECKED (T03$\sharp$ only: the conjunction of W3-T03ACC:T03S.1--T03S.4, clause map CHECKPOINT\_4.md §2; council R1, MANIFEST 9714aea4\ldots{}, astra=9.5 PASS, opus=9.4 PASS; LC-A, LC-B, LC-C VERIFIED). Not upgraded \ldots{}: T02(i$\sharp$) and T02(ii$\sharp$) PROVED, SF/XF VERIFIED; T01--T03 PROVED-LOCKED \ldots{}" · the card's alias \texttt{W3-POST:Thm-1@1.0.0} still reads "T03$\sharp$: PROVED; SF VERIFIED (W2-VT03R); XF VERIFIED (A-REV2)" \textbf{[FLAG H1]} & proved and independently verified (two independent verification passes) --- HQ ruling H1, 2026-11-16 \\ \addlinespace
Theorem 1, expectation display & 1.1 & \texttt{W3-POST2:C@1.0.5}: "PROVED-CHECKED, conditional on W3-POST:Prop-LB, [L1a$'$], W2-APP:POST, W3-POST:Thm-1"; \texttt{W3-POST:Prop-LB@1.0.0}: "PROVED (by this seat, unchecked)"; [L1a$'$] = \texttt{W3-POST/lemmas/L1a\_size\_tail.md} = \texttt{W3-POST:Lemma-1a'@1.0.0}: "PROVED (by this seat, unchecked)" (claims.tsv:343; corrected after the Sol read, D2-08); \texttt{W2-APP:POST@1.0.7}: "PROVED-CHECKED" & proved and checked, conditional on two supporting lemmas (a size-tail lemma and a lower bound on the posterior miss) that are proved but not yet independently checked \\ \addlinespace
Theorem 1$'$ & 1.1 & \texttt{W3-POST2:E@1.0.5}: "PROVED-CHECKED, conditional on W2-APP:POST, W3-POST:Thm-1"; \texttt{W3-POST2:A0@1.0.5}: "PROVED-CHECKED"; uses C case (iii) $\to$ inherits C's conditions & proved and checked, conditional on the same two supporting lemmas as the expectation display \\ \addlinespace
Remark (o) determinant growth & 1.1 & \texttt{W3-POST2:A1@1.0.5}: "PROVED-CHECKED, conditional on W3-POST:Thm-1, W2-APP:POST" & (remark) proved and checked \\ \addlinespace
Theorem 2 & 1.2 & \texttt{W3-POST:Thm-1} item 6 (status as row 1 / FLAG H1) + \texttt{W2-APP:POST@1.0.7}: "PROVED-CHECKED"; reader bridge B2.1 NEW & proved and independently checked; reader-form bridge (set-size arithmetic) awaits its final check \\ \addlinespace
Theorem 3 & 1.3 & \texttt{W3-PSHARP:A@1.0.0}, \texttt{W3-PSHARP:B@1.0.0}: "PROVED-CHECKED; council R2, MANIFEST 4513e05a\ldots{}, sol=9.2 PASS, opus=9.3 PASS; \ldots{}"; \texttt{W3-PSHARP-R:AE@1.0.0}, \texttt{W3-PSHARP-R:A2C@1.0.0}: "PROVED-CHECKED; council R2, MANIFEST fdfd4562\ldots{}, sol=9.5 PASS, opus=9.2 PASS"; exactness case B3.8 Sol-checked, upper half carries Thm 3.1's status & proved and independently checked \\ \addlinespace
Theorem 4 & 1.4 & \texttt{W3-MISR-R:CorGb-plus@1.0.1}: "PROVED-CHECKED, conditional on W3-POST2:C, W3-POST2:E, W3-POST2:A1, W3-POST2:A0, C-MIS:GDb, C-MIS:Prop-P; council R3, MANIFEST 052ea6df\ldots{}, sol=10.0 PASS, opus=9.6 PASS"; \texttt{C-MIS:Prop-P@1.0.0}: "R02 VERIFIED; R04 wording fixes checked in R05; council PASS"; \texttt{C-MIS:GD(b)@1.0.0}: "R02 VERIFIED-WITH-AMENDMENTS (applied); council PASS"; Example \texttt{W3-MISR-R:CorG-min@1.0.1}: "PROVED-CHECKED; council R3, \ldots{}"; remark \texttt{W3-MISR-R:CorH-plus@1.0.1}: "PROVED-CHECKED, conditional on W3-POST2:D; \ldots{}", \texttt{W3-MISR-R:CorH-U@1.0.1}: "PROVED-CHECKED, conditional on W3-POST2:D, W3-POST2:C, W3-POST2:A0, C-MIS:Prop-P; clause (G) also conditional on W3-POST2:E, W3-POST2:A1, C-MIS:GDb; \ldots{}"; headline masters \texttt{W3-MISR:MIS(i)@1.0.1}: "PROVED-CHECKED, conditional on W3-POST2:C; \ldots{}", MIS(i$'$), MIS(ii) "PROVED-CHECKED, conditional on W3-POST2:C, W3-POST2:E, W3-POST2:A1; \ldots{}" & proved and checked, conditional on the expectation form of Thm 3.1 (hence on its two supporting lemmas) and on two Gaussian-design lemmas that are verified and council-passed \\ \addlinespace
Theorem 9, (i)--(iii) & 2.1 & \texttt{W3-DIST:Thm-A@1.0.0}: "PROVED-CHECKED; council R2, MANIFEST 9719c817\ldots{} (= R2 c5cd9cdd re-frozen CITATION-ONLY as R2r), astra=9.3 PASS, opus=9.2 PASS" & proved and independently checked \\ \addlinespace
Theorem 9B & 2.1 & \texttt{W3-DIST:Thm-B@1.0.0}: "PROVED-CHECKED, conditional on W3-POST:Thm-1; council R2, \ldots{}" & proved and checked, conditional on Thm 3.1 (in-probability form), itself proved and checked (FLAG H1) \\ \addlinespace
Corollary 9C & 2.1 & \texttt{W3-DIST:Cor-C@1.0.0}: "PROVED-CHECKED, conditional on EX-A (exact arithmetic); council R2, \ldots{}" & proved and checked, under exact arithmetic \\ \addlinespace
Fixed-design remark (row 2.2) & 2.2 & \texttt{W3-POST2:D@1.1.5}: "PROVED-CHECKED, conditional on W2-APP:POST, W3-POST:Thm-1" & proved and checked (inputs checked) \\ \addlinespace
Theorem 5 & 3.3 & \texttt{W3-START2-R:START-R@1.0.0}, \texttt{W3-START2-R:BR@1.0.0}: "PROVED-CHECKED; council R2, MANIFEST 12bb52cd\ldots{} (= R2 90886d06 re-frozen CITATION-ONLY as R2r), sol=9.0 PASS, opus=9.3 PASS"; \texttt{W3-START:FLOOR@1.0.0}: "PROVED-CHECKED" & proved and independently checked \\ \addlinespace
Theorem 6, main & 3.1 & \texttt{W3-RM2:RM2-R@1.0.3}: "PROVED-CHECKED, conditional on W2-APP:POST; input versions W2-APP:POST@1.0.1; council R1, MANIFEST 00bbeab8\ldots{}, astra=9.5 PASS, opus=8.8 PASS; \ldots{}"; \texttt{W3-RM:RM-0@1.0.0}: "PROVED-CHECKED"; \texttt{W3-RM:E4@1.0.0}: "PROVED-CHECKED"; delivery \texttt{W2-UNI:F-UNI@2.1.0}: "PROVED-CHECKED"; reader bridges B6.3, B6.6 NEW; pending RM2-R swap \textbf{[FLAG H2]} & proved and checked (its one condition, the posterior transfer lemma, is itself proved and checked); reader-form bridges await their final check \\ \addlinespace
Theorem 7 & 3.2 & \texttt{W3-RM2:RM2-F@1.0.3}: "PROVED-CHECKED; council R1, MANIFEST 00bbeab8\ldots{}, astra=9.5 PASS, opus=8.8 PASS; \ldots{}"; RM-0, F-UNI as above; B7.4 NEW; pending RM2-R swap [FLAG H2] & proved and independently checked; reader-form delivery bridge awaits its final check \\ \addlinespace
Theorem 8 & 3.7 & \texttt{W3-SCR:A1@1.0.0}: "PROVED-CHECKED" & proved and independently checked \\ \addlinespace
Theorem 8$'$ & 3.7 & \texttt{W3-SCR:A-inf@1.0.0}: "PROVED-CHECKED"; \texttt{W3-SCR:SCR-A@1.0.0}: "PROVED-CHECKED" & proved and independently checked \\ \addlinespace
Theorem 10 & 3.6 & \texttt{W3-SAMP:C@1.1.4}: "PROVED-CHECKED"; start input \texttt{W3-START2-R:IF-START@1.0.0} (repin of \texttt{W3-START2:IF-START}, now SUPERSEDED-VALID): "PROVED-CHECKED; council R2, \ldots{}"; \texttt{W2-APP:POST@1.0.7}: "PROVED-CHECKED"; reader assembly B10.3 NEW & proved and independently checked (master); reader-form assembly awaits its final check \\ \addlinespace
Remark (row 3.5) per island & 3.5 & \texttt{W1-N1:Thm-N1@1.0.0}: "PROVED-CHECKED" & proved and independently checked \\ \addlinespace
Remark (row 3.5 / 4.1) pooled $\hat Z$ & 3.5, 4.1 & \texttt{W3-READER:PoolZ}: \textbf{no registry row on main} (card: "PROVED (by this seat, unchecked)") \textbf{[FLAG H3]} & proved (two-line argument), not yet independently checked \\ \addlinespace
Pooled PIP error $\le$ worst island & 4.1 & \texttt{W3-SAMP:LT\#@2.0.4}: "PROVED-CHECKED" & proved and independently checked \\ \addlinespace
Output-identical parallel engine & 4.2 & no registry row; engineering fact + note in progress & in progress (empirical; no theorem) \\ \addlinespace
Many islands / island weighting at 10$^6$ & 4.3 & no registry row; "OPEN" & open \\ \addlinespace
\bottomrule
\end{longtable}
\endgroup
\paragraph{Open flags for HQ.}
\begin{itemize}[leftmargin=1.5em]\footnotesize
\item \textbf{H1.} Label Theorem 1 (in-probability form) "proved and independently verified (two independent verification passes)", with a plain-words note on the two passes; never "council". Applied to Theorem 1, and to its dependants' notes (Theorem 2, Theorem 9B, row 2.2).
\item \textbf{H2.} "Revision pending" for rows 3.1/3.2 until the RM2 follow-up is recorded after the resume. Applied.
\item \textbf{H3.} Pooled evidence: "not yet independently checked". Applied.
\item \textbf{H4.} Prop-LB is not logged as implied by PSHARP:B (an unchecked, post-freeze item). The expectation form of Theorem 1, Theorem 1$'$ and Theorem 4 stay conditional on the lower-bound proposition and lemma [L1a$'$], not yet independently checked. Applied.
\end{itemize}
}{[Ledger table generated from STATUS\_LEDGER.md.]}

\end{document}
